% Source: DESIGN.md. Manually synchronized on 2026-08-20 because pandoc was unavailable.
% Options for packages loaded elsewhere
\PassOptionsToPackage{unicode}{hyperref}
\PassOptionsToPackage{hyphens}{url}
\documentclass[
  11pt,
]{article}
\usepackage{xcolor}
\usepackage[margin=1in]{geometry}
\usepackage{amsmath,amssymb}
\setcounter{secnumdepth}{-\maxdimen} % remove section numbering
\usepackage{iftex}
\ifPDFTeX
  \usepackage[T1]{fontenc}
  \usepackage[utf8]{inputenc}
  \usepackage{textcomp} % provide euro and other symbols
\else % if luatex or xetex
  \usepackage{unicode-math} % this also loads fontspec
  \defaultfontfeatures{Scale=MatchLowercase}
  \defaultfontfeatures[\rmfamily]{Ligatures=TeX,Scale=1}
\fi
\usepackage{lmodern}
\ifPDFTeX\else
  % xetex/luatex font selection
  % \setmonofont[]{Menlo}  % requires Menlo locally; Overleaf uses default Latin Modern Mono
\fi
% Use upquote if available, for straight quotes in verbatim environments
\IfFileExists{upquote.sty}{\usepackage{upquote}}{}
\IfFileExists{microtype.sty}{% use microtype if available
  \usepackage[]{microtype}
  \UseMicrotypeSet[protrusion]{basicmath} % disable protrusion for tt fonts
}{}
\makeatletter
\@ifundefined{KOMAClassName}{% if non-KOMA class
  \IfFileExists{parskip.sty}{%
    \usepackage{parskip}
  }{% else
    \setlength{\parindent}{0pt}
    \setlength{\parskip}{6pt plus 2pt minus 1pt}}
}{% if KOMA class
  \KOMAoptions{parskip=half}}
\makeatother
\usepackage{longtable,booktabs,array}
\newcounter{none} % for unnumbered tables
\usepackage{calc} % for calculating minipage widths
% Correct order of tables after \paragraph or \subparagraph
\usepackage{etoolbox}
\makeatletter
\patchcmd\longtable{\par}{\if@noskipsec\mbox{}\fi\par}{}{}
\makeatother
% Allow footnotes in longtable head/foot
\IfFileExists{footnotehyper.sty}{\usepackage{footnotehyper}}{\usepackage{footnote}}
\makesavenoteenv{longtable}
\setlength{\emergencystretch}{3em} % prevent overfull lines
\providecommand{\tightlist}{%
  \setlength{\itemsep}{0pt}\setlength{\parskip}{0pt}}
\usepackage{amsmath,amssymb,amsthm}
\usepackage{mathtools}
\def\C{\mathbb{C}}
\def\R{\mathbb{R}}
\def\Z{\mathbb{Z}}
\def\N{\mathbb{N}}
\def\Q{\mathbb{Q}}
\def\H{\mathcal{H}}
\def\B{\mathcal{B}}
\def\K{\mathcal{K}}
\def\A{\mathcal{A}}
\def\D{D}
\DeclareMathOperator{\ind}{ind}
\DeclareMathOperator{\Tr}{Tr}
\DeclareMathOperator{\dom}{dom}
\usepackage{bookmark}
\IfFileExists{xurl.sty}{\usepackage{xurl}}{} % add URL line breaks if available
\urlstyle{same}
\hypersetup{
  pdftitle={Spectral Triples and the Index Pairing},
  pdfauthor={Project notes},
  hidelinks,
  pdfcreator={LaTeX via pandoc}}

\title{Spectral Triples and the Index Pairing}
\usepackage{etoolbox}
\makeatletter
\providecommand{\subtitle}[1]{% add subtitle to \maketitle
  \apptocmd{\@title}{\par {\large #1 \par}}{}{}
}
\makeatother
\subtitle{Formalization design --- see PLAN.md for implementation status}
\author{Project notes}
\date{2026-08-20}

\begin{document}
\maketitle

\section{1. Goal}\label{goal}

\texttt{DESIGN.md} is the source document for this roadmap;
\texttt{DESIGN.tex} is its secondary presentation artifact. Status claims
name a Lean declaration when they describe a formalized result. Unqualified
index-theorem statements are classical targets, not claims about the present
codebase.

Formalize, in Lean 4 with Mathlib, the following triad of objects and
the relations between them:

\begin{enumerate}
\def\labelenumi{(\alph{enumi})}
\tightlist
\item
  the abstract notion of a \textbf{\(Z_2\)-graded spectral triple}
  \((\mathcal{A}, \mathcal{H}, D, \gamma)\),
\item
  the \textbf{Fredholm index pairing}
  \(\bigl\langle [D],\,[p]\bigr\rangle = \mathop{\mathrm{ind}}\bigl(p D^+ p\bigr) \in \mathbb{Z}\)
  of an even spectral triple with a projection \(p \in \mathcal{A}\),
\item
  the \textbf{canonical spectral triple of a closed Riemannian spin
  manifold} \((M, g)\), i.e.~the triple
  \(\bigl(C^{\infty}(M),\, L^2(S),\, D_M,\, \gamma_M\bigr)\) built from
  the spinor bundle \(S \to M\) and the Dirac operator \(D_M\).
\end{enumerate}

The mathematical content is classical (Connes 1994; GBF 2001;
Higson--Roe 2000). The intended contribution is a Lean encoding and an
eventual \emph{machine-checked} verification that the manifold
construction satisfies the abstract axioms. The current implementation
has reached the abstract core and Fourier models, not the general
manifold construction.

\subsection{\texorpdfstring{1.1 Motivation: why both an algebra
\(\mathcal{A}\) and a Hilbert space
\(\mathcal{H}\)?}{1.1 Motivation: why both an algebra \textbackslash mathcal\{A\} and a Hilbert space \textbackslash mathcal\{H\}?}}\label{motivation-why-both-an-algebra-mathcala-and-a-hilbert-space-mathcalh}

The bipartite data \((\mathcal{A}, \mathcal{H}, D)\) encodes a
deliberate split: \(\mathcal{A}\) carries the \emph{topology} (what the
space looks like), and \(\mathcal{H}\) together with \(D\) carries the
\emph{calculus} (how to differentiate and measure on it). Each half
alone is too featureless to recover geometry; the spectral triple is
their collision.

\textbf{Algebra alone is too rubbery.} By Gelfand--Naimark a commutative
\(C^*\)-algebra fully determines a (locally compact Hausdorff)
topological space, but topology cannot see distances, angles, or
volumes: \(C(X)\) does not distinguish a unit sphere from a sphere of
radius 100. And on a generic noncommutative algebra there are no
nontrivial outer derivations at all, so the very notion of
``differentiating an element of \(\mathcal{A}\)'' is absent from
\(\mathcal{A}\) alone.

\textbf{Hilbert space alone is amnesiac.} Every separable
infinite-dimensional Hilbert space is unitarily isomorphic to
\(\ell^2(\mathbb{N})\): \(L^2(S^2)\) and \(L^2(T^2)\) are
indistinguishable as Hilbert spaces. Adding a self-adjoint operator
\(D\) does not fix this --- by Milnor (1964) and the negative answer to
Kac's ``Can one hear the shape of a drum?'' (Gordon--Webb--Wolpert
1992), isospectral non-isometric manifolds exist. The spectrum of \(D\)
alone misses the manifold.

\textbf{The representation \(\pi\) localises.} A \(*\)-representation
\(\pi\colon \mathcal{A}\to \mathcal{B}(\mathcal{H})\) breaks the
universality of \(\mathcal{H}\) by anchoring it to the topological shape
recorded by \(\mathcal{A}\). For the canonical Dirac triple,
\((\pi(f)\psi)(x) = f(x)\,\psi(x)\) is pointwise multiplication --- this
is what lets a vector in \(\mathcal{H}\) know \emph{which point of \(M\)
it lives over}. Connes' reconstruction theorem (1996; 2008 final form)
makes this precise: from \((\mathcal{A}, \mathcal{H}, D, \gamma)\)
satisfying suitable additional axioms, the smooth Riemannian spin
manifold \(M\) is recovered.

\textbf{The commutator is the derivative.} \(D\) acts on
\(\mathcal{H}\), not on \(\mathcal{A}\), so \(a \in \mathcal{A}\) cannot
be differentiated by \(D\) directly. But both \(D\) and \(\pi(a)\) live
in \(\mathcal{B}(\mathcal{H})\), so they can be multiplied; their
commutator \[
[D, \pi(a)]\,\psi \;=\; D(a\psi) - a\,(D\psi) \;=\; c(da)\,\psi
\] \emph{is} Clifford multiplication by \(df\) in the canonical manifold
case. Boundedness of \([D, \pi(a)]\) (axiom (3) of §2.1) is the abstract
statement that ``\(a\) is Lipschitz.'' The derivative is born from the
failure of \(\mathcal{A}\) to commute with \(D\) --- and that failure
exists only because both act on a common third object \(\mathcal{H}\).

\textbf{The pairing: \(\mathcal{A}\) asks, \(\mathcal{H}+ D\) answers.}
A projection \(p \in \mathcal{A}\) defines a K-theory class --- a
topological question (``does this bundle have holes?''). The Fredholm
index of \(p D^+ p \colon p\mathcal{H}^+ \to p\mathcal{H}^-\) is the
analytic answer \(\langle [D], [p]\rangle \in \mathbb{Z}\). Neither half
produces the integer alone: \(\mathcal{A}\) supplies the topological
lock, \(\mathcal{H}+ D\) supplies the analytic key. This is the classical
index pairing of §2.2. Current Phase-2 code contains a generic bounded
Fredholm index and a separate \texttt{gradedKernelIndex}; the
projection-compressed chiral pairing remains a roadmap item.

\textbf{Heisenberg's legacy.} Connes modelled the triple on the
Heisenberg picture of quantum mechanics: \(\mathcal{A}=\) observables,
\(\mathcal{H}=\) states, \(D =\) Hamiltonian. Just as QM needs all three
together to mean anything, ``space'' in noncommutative geometry is born
from the interaction of \(\mathcal{A}\) and \(D\) on a common
\(\mathcal{H}\).

\section{2. Mathematical content}\label{mathematical-content}

\subsection{\texorpdfstring{2.1 Abstract \(Z_2\)-graded spectral
triple}{2.1 Abstract Z\_2-graded spectral triple}}\label{abstract-z_2-graded-spectral-triple}

Let \(\mathcal{A}\) be a unital \(*\)-algebra over \(\mathbb{C}\) and
\(\mathcal{H}\) a separable complex Hilbert space.

\textbf{Definition (Connes).} A \emph{spectral triple}
\((\mathcal{A}, \mathcal{H}, D)\) consists of:

\begin{enumerate}
\def\labelenumi{\arabic{enumi}.}
\tightlist
\item
  a \(*\)-representation
  \(\pi \colon \mathcal{A}\to \mathcal{B}(\mathcal{H})\) by bounded
  operators;
\item
  a self-adjoint operator \(D\) on \(\mathcal{H}\) with dense domain
  \(\mathop{\mathrm{dom}}D\), such that
  \((D+ i)^{-1} \in \mathcal{K}(\mathcal{H})\) (compact resolvent);
\item
  for every \(a \in \mathcal{A}\),
  \(\pi(a)\,\mathop{\mathrm{dom}}D\subseteq \mathop{\mathrm{dom}}D\) and
  the commutator
  \([D, \pi(a)]\colon \mathop{\mathrm{dom}}D\to \mathcal{H}\) extends to
  a bounded operator on \(\mathcal{H}\).
\end{enumerate}

The triple is \textbf{\(Z_2\)-graded} (or \emph{even}) if there exists a
bounded self-adjoint involution
\(\gamma\colon \mathcal{H}\to \mathcal{H}\) (i.e.~\(\gamma^* = \gamma\)
and \(\gamma^2 = 1\)) with

\[
[\gamma,\pi(a)] = 0 \quad \forall\,a \in \mathcal{A},
\qquad
\gamma\,D+ D\,\gamma = 0.
\]

The grading splits \(\mathcal{H}= \mathcal{H}^+ \oplus \mathcal{H}^-\)
into the \(\pm 1\) eigenspaces of \(\gamma\). The representation \(\pi\)
preserves each summand; \(D\) exchanges them. Writing
\(D= \begin{pmatrix} 0 & D^- \\ D^+ & 0 \end{pmatrix}\) with respect to
the splitting, self-adjointness of \(D\) gives \(D^- = (D^+)^*\).

\subsection{2.2 Index pairing}\label{index-pairing}

Let \((\mathcal{A}, \mathcal{H}, D, \gamma)\) be a \(Z_2\)-graded
spectral triple and let \(p = p^* = p^2 \in \mathcal{A}\) be a
projection.

\textbf{Classical theorem (Connes; formalization target).} The operator
\(p D^+ p \colon p\mathcal{H}^+ \to p\mathcal{H}^-\) is Fredholm, and
its index depends only on the class \([p] \in K_0(\mathcal{A})\)
(K-theory of the algebra) and on the class \([D] \in K^0(\mathcal{A})\)
(even analytic K-homology). The pairing \[
\bigl\langle [D],\,[p]\bigr\rangle
\;:=\; \mathop{\mathrm{ind}}\bigl(p D^+ p\bigr)
\;=\; \dim\ker\bigl(p D^+ p\bigr) \;-\; \dim\operatorname{coker}\bigl(p D^+ p\bigr)
\;\in\; \mathbb{Z}
\] is bilinear and additive. Self-adjointness of \(D\) identifies the
cokernel with \(\operatorname{coker}(p D^+ p) \cong \ker(p D^- p)\) as a
subspace of \(p\mathcal{H}^-\), since \((pD^+ p)^* = p D^- p\) under the
splitting \(\mathcal{H}= \mathcal{H}^+ \oplus \mathcal{H}^-\) (see
§2.1).

Fredholmness reduces to two facts:

\begin{itemize}
\tightlist
\item
  \(p D^+ p\) has a parametrix modulo compacts, built from
  \((D+ i)^{-1}\) (compact by the resolvent axiom) and the bounded
  commutator \([D, \pi(p)]\) (bounded by the commutator axiom);
\item
  compact resolvent + bounded commutator \(\Rightarrow\) compactness of
  the remainder term, so Atkinson's theorem applies.
\end{itemize}

\textbf{Bounded form.} Replace \(D\) by \(F := D\,(1 + D^2)^{-1/2}\)
(the \emph{bounded transform} of \(D\)). Then \(F\) is bounded
self-adjoint, \(F^2 - 1\) and \([F, \pi(a)]\) are compact for every
\(a \in \mathcal{A}\), and the pairing reads \[
\bigl\langle [F],\,[p]\bigr\rangle
\;=\; \mathop{\mathrm{ind}}\bigl(p F^+ p\colon p\mathcal{H}^+ \to p\mathcal{H}^-\bigr) \in \mathbb{Z}.
\] Here \(F^+\colon \mathcal{H}^+ \to \mathcal{H}^-\) is the
off-diagonal block of \(F\) under the \(\gamma\)-grading. The two
pictures give the same integer (GBF §9.4).

\textbf{Present Lean boundary.} \texttt{Fredholm.lean} defines the
Fredholm predicate and index for bounded linear maps, while
\texttt{Index.lean} defines \texttt{gradedKernelIndex} for the full
self-adjoint unbounded operator and proves its kernel finite-dimensional
under compact resolvent. The chiral restriction \(D^+\), its Fredholmness,
projection compression, and equality with \texttt{gradedKernelIndex} are
not yet formalized.

\subsection{2.3 Canonical spectral triple of a Riemannian spin
manifold}\label{canonical-spectral-triple-of-a-riemannian-spin-manifold}

Let \((M, g)\) be a closed Riemannian manifold of dimension \(n\),
equipped with a spin structure. The construction proceeds in stages.

\textbf{(i) Clifford algebra and spinor bundle.} The cotangent bundle
\(T^*M\) together with \(g\) gives a Clifford bundle
\(\mathrm{Cl}(T^*M)\) whose fibre \(\mathrm{Cl}_n\) is generated by
\(T^*_x M\) with relations \(v \cdot w + w \cdot v = -2 g(v, w)\). A
\textbf{spin structure} on \(M\) is a principal
\(\mathrm{Spin}(n)\)-bundle \(P_{\mathrm{Spin}} \to M\) together with a
double cover \(P_{\mathrm{Spin}} \to P_{\mathrm{SO}}\) of the oriented
orthonormal frame bundle, equivariant for the standard double cover
\(\mathrm{Spin}(n) \to \mathrm{SO}(n)\). Such a lift exists iff the
second Stiefel--Whitney class \(w_2(M) \in H^2(M; \mathbb{Z}/2)\)
vanishes; when it does, the set of spin structures is a torsor over
\(H^1(M; \mathbb{Z}/2)\).

Associated to the standard complex spinor representation
\(\sigma\colon \mathrm{Cl}_n \otimes_{\mathbb{R}} \mathbb{C}\to \mathrm{End}(\Sigma_n)\)
(irreducible of dimension \(2^{\lfloor n/2 \rfloor}\); unique for \(n\)
even, two inequivalent choices distinguished by
\(\sigma(\omega) = \pm 1\) for \(n\) odd), one obtains the
\textbf{spinor bundle} \[
S \;:=\; P_{\mathrm{Spin}} \times_{\mathrm{Spin}(n)} \Sigma_n
\;\longrightarrow\; M.
\] Pointwise Clifford multiplication \(c\colon T^*M \otimes S \to S\) is
fibrewise the action \(\sigma\).

\textbf{(ii) Hilbert space.} With the Riemannian volume form, define
\(\mathcal{H}:= L^2(M, S)\), the completion of smooth compactly
supported sections \(C_c^{\infty}(M, S)\) under
\(\langle \psi, \varphi\rangle := \int_M \langle \psi(x), \varphi(x)\rangle_{S_x}\,d\mathrm{vol}_g(x)\).

\textbf{(iii) Algebra and representation.} Let
\(\mathcal{A}:= C^{\infty}(M)\) act on \(\mathcal{H}\) by pointwise
multiplication: \((\pi(f)\psi)(x) := f(x)\,\psi(x)\).

\textbf{(iv) Dirac operator.} The Levi--Civita connection on \(TM\)
lifts to a connection \(\nabla^S\) on \(S\). The Dirac operator is \[
D_M \;:=\; c \circ \nabla^S
\;\colon\; \Gamma(S) \to \Gamma(S),
\] in local coordinates \(D_M = \sum_i c(e^i)\,\nabla^S_{e_i}\) for any
local orthonormal frame \((e_i)\) with dual coframe \((e^i)\).

\textbf{(v) Grading (when \(n\) is even).} In even dimension the spinor
representation splits as \(\Sigma_n = \Sigma_n^+ \oplus \Sigma_n^-\)
under the \textbf{chirality element} \[
\gamma_M \;:=\; i^{n/2}\, c(e^1) c(e^2) \cdots c(e^n),
\] which is self-adjoint, squares to \(1\), commutes with even Clifford
elements, anti-commutes with odd ones. Hence \(\gamma_M\) commutes with
\(\pi(f)\) (a function is a degree-\(0\) object) and anti-commutes with
\(D_M\) (which is degree-\(1\)).

\textbf{Theorem (canonical triple).} With the above data,
\((C^{\infty}(M), L^2(M,S), D_M, \gamma_M)\) is a \(Z_2\)-graded
spectral triple. Specifically (writing \(D_M\) for both the symmetric
operator on \(C_c^\infty(M, S)\) and its closure):

\begin{itemize}
\tightlist
\item
  \(D_M\) is essentially self-adjoint on \(C_c^{\infty}(M,S)\) ---
  closed \(M\) is geodesically complete, so Chernoff's theorem applies.
  We identify \(D_M\) with its self-adjoint closure throughout.
\item
  \((D_M + i)^{-1}\) is compact, by Rellich--Kondrachov: the inclusion
  \(H^1(M,S) \hookrightarrow L^2(M,S)\) is compact for compact \(M\),
  and \((D_M + i)^{-1}\) factors through this inclusion since
  \((D_M + i)^{-1}\colon L^2 \to H^1\) is bounded by elliptic
  regularity.
\item
  For \(f \in C^{\infty}(M)\), the commutator
  \([D_M, \pi(f)]\colon C_c^\infty(M, S) \to C_c^\infty(M, S)\) equals
  \(c(df)\), pointwise Clifford multiplication by the (bounded smooth)
  one-form \(df\). It extends by boundedness to all of \(L^2(M, S)\)
  with operator norm \(\|df\|_{\infty}\) (the sup norm of \(df\) in the
  metric induced by \(g\)).
\item
  \(\gamma_M\) commutes with \(\pi(C^{\infty}(M))\) (functions are
  degree-0 Clifford elements) and anti-commutes with \(D_M\) (which is
  degree-1).
\end{itemize}

Atiyah--Singer specializes the index pairing: for a Hermitian vector
bundle \(E \to M\) with projection \(p_E\) encoding \(E\), the index
\(\langle [D_M], [p_E] \rangle = \mathop{\mathrm{ind}}(D_E^+)\) where
\(D_E\) is the twisted Dirac operator, equal to the topological index of
\(E\) (Â-genus times Chern character).

\subsection{2.4 Generalizations: equivariance and
families}\label{generalizations-equivariance-and-families}

The framework admits two important extensions, both well established in
the literature and both fitting naturally on top of the abstract
\(Z_2\)-graded spectral triple of §2.1.

\textbf{Equivariance.} Let \(G\) be a locally compact group acting on
\(\mathcal{A}\) by \(*\)-automorphisms
\(\sigma\colon G \to \mathrm{Aut}(\mathcal{A})\). A
\emph{\(G\)-equivariant} (graded) spectral triple adds a
strongly-continuous unitary representation
\(U\colon G \to \mathcal{U}(\mathcal{H})\) satisfying the covariance
condition \[
U(g)\,\pi(a)\,U(g)^* = \pi\!\bigl(\sigma_g(a)\bigr)
\qquad \forall\, g \in G,\ a \in \mathcal{A},
\] together with \(U(g)D= DU(g)\) and (in the graded case)
\(U(g)\gamma = \gamma U(g)\) for every \(g\). For \textbf{compact} \(G\)
the index pairing refines to a homomorphism
\(K^G_0(\mathcal{A}) \to R(G)\) into the representation ring of \(G\);
specializing at the trivial representation recovers the integer index.
For non-compact \(G\) one works with the \(K\)-theory of the reduced
group \(C^*\)-algebra \(K_0(C^*_r(G))\) in place of \(R(G)\), via the
Baum--Connes assembly map; this generality is out of scope for the
initial formalization.

Prototype examples: Atiyah--Singer for \(G\)-equivariant elliptic
operators on a \(G\)-manifold (with isometric \(G\)-action and
\(G\)-invariant spin structure); Connes--Moscovici's transverse
fundamental class on a foliation; spectral triples for quantum groups.

\textbf{Families.} Replace the Hilbert space \(\mathcal{H}\) with a
\textbf{Hilbert \(C^*\)-module} \(\mathcal{E}\) over a \(C^*\)-algebra
\(B\) (think \(B = C(X)\) for a compact Hausdorff base \(X\), so that
\(\mathcal{E}\) is the section space of a continuous field of Hilbert
spaces over \(X\)). The representation \(\pi\) now takes values in the
adjointable \(B\)-linear operators \(\mathcal{L}_B(\mathcal{E})\), and
``compact resolvent'' means
\((D+ i)^{-1} \in \mathcal{K}_B(\mathcal{E})\), the \(B\)-compact
operators (the closed two-sided ideal generated by the rank-one
\(B\)-linear maps
\(\theta_{\xi,\eta}(\zeta) := \xi\,\langle \eta, \zeta\rangle_B\)).

This is precisely an \textbf{unbounded Kasparov
\((\mathcal{A}, B)\)-bimodule}; the bounded form (replace \(D\) by
\(F\)) is Kasparov's original definition. The index lifts to a class in
\(K_0(B)\); for \(B = C(X)\) this recovers the Atiyah--Singer family
index theorem. The whole programme sits inside \textbf{Kasparov's
bivariant K-theory} \(KK(\mathcal{A}, B)\), and our current setting is
the special case \(B = \mathbb{C}\) (where
\(K_0(\mathbb{C}) = \mathbb{Z}\)).

The two extensions combine: \(G\)-equivariant families give classes in
\(KK^G(\mathcal{A}, B)\), encompassing equivariant family index
theorems.

\subsection{2.5 Super-Bakry--Émery
curvature}\label{super-bakryuxe9mery-curvature}

A spectral triple comes with enough structure to support a synthetic
\textbf{Ricci-curvature lower bound} in the sense of Bakry--Émery,
adapted to the Z₂-graded setting. The generator is \(L := D^2\)
(positive, self-adjoint on \(\mathcal{H}\)); the ``gradient'' is the
noncommutative derivative \(\delta(a) := [D, a]\), an odd derivation
\(\mathcal{A}\to \mathcal{B}(\mathcal{H})\).

\textbf{Carré du champ.} Define \[
\Gamma(a, b) \;:=\; \tfrac{1}{2}\bigl(\delta(a)^*\,\delta(b) + \delta(b^*)\,\delta(a^*)^*\bigr)
\qquad (a, b \in \mathcal{A}),
\] a sesquilinear form valued in \(\mathcal{B}(\mathcal{H})\)
(Cipriani--Sauvageot; agrees with
\(\Gamma(a, b) = \tfrac{1}{2}(L(a^*b) - a^*\,Lb - (La^*)\,b)\) when
\(\mathcal{A}\) acts as multiplication operators and \(L\) is a Markov
generator). The \textbf{iterated carré du champ} is \[
\Gamma_2(a, b) \;:=\; \tfrac{1}{2}\bigl(L\,\Gamma(a, b) - \Gamma(a, Lb) - \Gamma(La, b)\bigr).
\] For real \(\mathcal{A}\) and self-adjoint \(a\),
\(\Gamma(a, a) \geq 0\) in the operator order.

\textbf{Super-\(CD(\rho, \infty)\) condition.} The spectral triple
satisfies super-Bakry--Émery curvature \(\geq \rho \in \mathbb{R}\) if
\[
\Gamma_2(a, a) \;\geq\; \rho\,\Gamma(a, a) \qquad \forall\,a \in \mathcal{A}
\] in the operator-ordering sense (positive operator inequality on
\(\mathcal{H}\)). The ``super'' qualifier marks two non-classical
features: the generator is \(L = D^2\) with \(D\) odd (Dirac-type), and
the underlying space is Z₂-graded by \(\gamma\).

\textbf{Equivalent gradient-decay form} (used in \texttt{graphops-qft}).
Let \(P_t := e^{-tL}\) be the heat semigroup. The \(CD(\rho, \infty)\)
inequality is equivalent to \[
\Gamma(P_t a,\,P_t a) \;\leq\; e^{-2\rho t}\,P_t\bigl(\Gamma(a, a)\bigr)
\qquad \forall\, t \geq 0,\ a \in \mathcal{A}.
\] Equivalently in expectation:
\(\|\nabla P_t a\|^2 \leq e^{-2\rho t}\,P_t \|\nabla a\|^2\) where
\(\|\nabla a\|^2 := \Gamma(a, a)\) understood as an
\(\mathcal{H}\)-bilinear form.

\textbf{Specialization to the canonical Dirac triple.} Via the
\textbf{Lichnerowicz formula} \[
D^2 \;=\; \nabla^{S\,*}\nabla^S \,+\, \tfrac{1}{4} R_{\mathrm{scal}}(g),
\] combined with the Bochner formula for the rough Laplacian, the
Bakry--Émery \(CD\) inequality for the canonical Dirac spectral triple
of a closed Riemannian spin manifold \((M, g)\) reduces to a pointwise
lower bound on \(\mathrm{Ric}_g\). Concretely: \(\mathrm{Ric}_g \geq K\)
pointwise implies super-\(CD(K/4, \infty)\) for the spectral triple. The
factor \(1/4\) is the Lichnerowicz coefficient.

\textbf{Bridge to \texttt{graphops-qft}.} The \texttt{SusyGraphop}
\texttt{(H₊,\ H₋,\ Q)} of \texttt{graphops-qft} is exactly a
finite-dimensional Fredholm module with \(Q\) playing the role of \(F\)
(legal since \(F = Q\) in finite dimension --- the bounded transform is
the identity up to scaling). The gradient \(\nabla = Q\) (acting
\(C^0 \to C^1\) on a graph) is the supercharge; the super-BE curvature
defined there agrees with the carré-du-champ form above with \(L = Q^2\)
on the even side. So the spectral-triples project contains
\texttt{graphops-qft}'s finite case as an instance of a graded Fredholm
module; integrating super-BE here gives both frameworks a common
abstraction.

\subsection{2.6 Other instances of the
framework}\label{other-instances-of-the-framework}

\emph{Extra, not core.} The abstract data of an even graded Fredholm
module
\((\mathcal{A}, \mathcal{H}= \mathcal{H}^+ \oplus \mathcal{H}^-, F, \gamma)\)
--- equivalently a pair of Hilbert spaces with an odd operator \(Q\)
between them --- has several classical instantiations beyond the Dirac
case of §2.3. Each comes with its own deformation parameter or extra
structure; in each, the index pairing reduces to a familiar topological
invariant.

\textbf{Hodge--de Rham (the signature/Euler operator).} Take
\(\mathcal{H}= L^2(\Lambda^* M)\) (square-integrable forms of all
degrees) on a closed oriented Riemannian manifold,
\(\gamma = (-1)^{\deg}\) for the Euler operator (or the Hodge-\(*\)
chirality for the signature operator), and \(D = d + d^*\). The Witten
index \(\dim\ker D|_{\mathrm{even}} -
\dim\ker D|_{\mathrm{odd}}\) recovers the Euler characteristic
\(\chi(M) = \sum (-1)^k \beta_k\) (resp. the signature \(\sigma(M)\) in
the oriented even-dim chirality case).

\textbf{Witten--Morse deformation.} A Morse function
\(f \colon M \to \mathbb{R}\) defines a one-parameter family
\(d_t \;:=\; e^{-tf}\,d\,e^{tf} \;=\; d \,+\, t\,df \wedge\) of
deformations of \(d\), hence a family of even graded Fredholm modules
\(\bigl(C^\infty(M),\, L^2(\Lambda^*M),\, d_t + d_t^*,\, (-1)^{\deg}\bigr)\)
for \(t \in [0, \infty)\). The Witten Laplacian
\(\Delta_t := (d_t + d_t^*)^2\) agrees with the standard Hodge Laplacian
at \(t = 0\) and develops a semiclassical large-\(t\) asymptotic of the
form
\(\Delta_t \;\sim\; \Delta \,+\, t^2 |\nabla f|^2 \,+\, t\,(\mathcal{L}_{\nabla f} + \mathcal{L}_{\nabla f}^*)\)
in which the potential \(t^2 |\nabla f|^2\) confines low-energy
eigenfunctions to a shrinking neighbourhood of the critical points of
\(f\) (Witten 1982; rigorous via Helffer--Sjöstrand 1985). The Witten
index is \(t\)-independent and equals \(\chi(M)\); the count of
low-energy modes by Morse index gives the \textbf{Morse inequalities};
the exponentially-small eigenvalues at large \(t\) encode the
\textbf{Morse complex differential} (counted gradient flow lines, à la
Floer/Bismut).

\textbf{Morse-complex limit as a finite \texttt{SusyGraphop}.} In the
\(t \to \infty\) limit the smooth Witten triple ``collapses'' onto a
finite-dimensional chain complex: one generator per critical point of
\(f\), graded by Morse index, with differential given by flow-line
counts. This finite complex \textbf{is} a \texttt{SusyGraphop} in the
sense of \texttt{graphops-qft}. So Witten--Morse provides a
\emph{smooth} discrete-continuum bridge --- complementing the
graph-approximation bridge --- in which the same K-homology class is
represented by the smooth spectral triple (\(t = 0\)) and by the finite
Morse complex (\(t = \infty\)).

\textbf{Manifolds with boundary --- deliberately out of scope.} On a
Riemannian manifold \(M\) with \(\partial M \neq \emptyset\), the Dirac
operator is not essentially self-adjoint, requiring a choice of
self-adjoint extension (APS, MIT-bag, Riemannian/Dirichlet, \ldots) ---
each giving a different operator with a different index. The
Atiyah--Patodi--Singer (1975) formula
\(\mathrm{ind}(D_{\mathrm{APS}}^+) = \int_M \widehat{A}(M)\,\mathrm{ch}(E) - \tfrac{1}{2}(\eta(D_{\partial M}) + h(\partial M))\)
introduces the non-local \textbf{eta-invariant}
\(\eta(s) := \sum_{\lambda \neq 0} \mathrm{sgn}(\lambda)\,|\lambda|^{-s}\)
of the boundary Dirac operator --- a real (non-integer) spectral
asymmetry, requiring meromorphic continuation in \(s\) and well outside
our K-theoretic / Fredholm setup. \emph{In the bounded picture, the
operator-algebraic signature of the boundary obstruction is precisely
the failure of \(F^2 - 1\) to be compact:} the boundary algebra
\(C(\partial M)\) emerges as the cokernel modulo compacts (this is the
relative-K-homology setting of Forsyth--Mesland--Rennie 2019,
Goffeng--Mesland 2015, and the relative-K-homology / ``spectral triples
with boundary'' literature). We therefore \emph{deliberately exclude}
manifolds with boundary from this project --- the classical bounded
Fredholm-module condition \(F^2-1\in\mathcal{K}(\mathcal{H})\) enforces it. The
closed-manifold case is enough to formalize the index pairing and the
canonical construction; the APS framework is a separate project of
comparable scope.

\textbf{Hodge--Dolbeault.} For a complex manifold \((X, J)\),
\(D = \overline{\partial} + \overline{\partial}^*\) on
\(\mathcal{H}= L^2(\Lambda^{0,*})\) with \(\gamma = (-1)^{q}\) on
\(\Lambda^{0,q}\) gives the Dolbeault spectral triple. Its index is the
holomorphic Euler characteristic
\(\chi(\mathcal{O}_X) = \sum (-1)^q h^{0,q}\); twisting by a holomorphic
vector bundle \(E\) gives \(\chi(E)\) via Hirzebruch--Riemann--Roch.
This is the framework specialization that underwrites our Phase 3.5
\(T^2\) line-bundle test (where it reduces to Riemann--Roch on a genus-1
curve).

\subsection{2.7 Connes' spectral
distance}\label{connes-spectral-distance}

The data of a spectral triple supplies not only an \emph{integer-valued}
invariant (the index pairing of §2.2) but also a \emph{metric}. The
mechanism: the bounded commutator \([D, \pi(a)]\) is the noncommutative
gradient of \(a\), and a noncommutative analogue of the classical
formula \(d(x, y) = \sup\{|f(x) - f(y)| : \|\nabla f\|_\infty \leq 1\}\)
yields a distance on the state space of \(\mathcal{A}\).

\textbf{Definition (Connes).} Let \((\mathcal{A}, \mathcal{H}, D)\) be a
spectral triple and let \(\mathcal{S}(\mathcal{A})\) denote its state
space (positive linear functionals
\(\phi\colon \mathcal{A}\to \mathbb{C}\) of norm 1). The
\textbf{spectral distance} is \[
d_D(\phi, \psi) \;:=\; \sup\Bigl\{\,\bigl|\phi(a) - \psi(a)\bigr| \;:\; a \in \mathcal{A},\ a = a^*,\ \bigl\|[D, \pi(a)]\bigr\|_{\mathcal{B}(\mathcal{H})} \leq 1\,\Bigr\}
\qquad (\phi, \psi \in \mathcal{S}(\mathcal{A})).
\] It takes values in \([0, +\infty]\), defines an extended metric on
\(\mathcal{S}(\mathcal{A})\), and depends only on the gauge-equivalence
class of \(D\).

Three specializations:

\begin{itemize}
\tightlist
\item
  \textbf{Manifold (commutative case).} For the canonical Dirac triple
  of a closed Riemannian spin manifold \((M, g)\), the algebra is
  \(\mathcal{A}= C^\infty(M)\) and \emph{pure states} are evaluation
  functionals \(\delta_x(f) := f(x)\) for \(x \in M\). Connes' theorem:
  \[d_D(\delta_x, \delta_y) \;=\; d_g(x, y) \qquad \forall\, x, y \in M,\]
  the Riemannian geodesic distance. This is the foundational result
  showing that spectral data \(\bigl(C^\infty(M), L^2(S), D_M\bigr)\)
  recovers the metric of \((M, g)\) --- a key step toward Connes'
  reconstruction theorem.
\item
  \textbf{Graph (commutative finite case).} For a finite graph viewed as
  a finite spectral triple (\texttt{graphops-qft}
  \texttt{SpectralTriple}, with \(\mathcal{A}\) the algebra of functions
  on vertices and \(D = d + d^*\)), the spectral distance between vertex
  evaluation states equals the graph (shortest-path) distance. The
  condition \(\|[D, f]\| \leq 1\) becomes \(|f(x) - f(y)| \leq 1\) along
  edges --- exactly the 1-Lipschitz condition for the combinatorial
  metric.
\item
  \textbf{Noncommutative case (generalized Wasserstein).} For a
  genuinely noncommutative \(\mathcal{A}\) --- e.g., matrix algebras,
  group \(C^*\)-algebras of discrete groups, fuzzy spheres --- the
  formula yields the \textbf{Connes--Rieffel metric} on
  \(\mathcal{S}(\mathcal{A})\), an analogue of the Wasserstein-1
  (earth-mover) distance. The Lipschitz seminorm
  \(L(a) := \|[D, \pi(a)]\|\) plays the role of a Lipschitz constant,
  and \(d_D\) is the dual transport distance for \(L\). Rieffel (1999)
  identified the conditions on \(L\) under which \(d_D\) metrizes the
  weak-\(*\) topology on \(\mathcal{S}(\mathcal{A})\) --- a ``compact
  quantum metric space'' structure.
\end{itemize}

\textbf{Relation to the index pairing.} The two invariants extracted
from spectral data are complementary:

{\def\LTcaptype{none} % do not increment counter
\begin{longtable}[]{@{}
  >{\raggedright\arraybackslash}p{(\linewidth - 4\tabcolsep) * \real{0.2500}}
  >{\raggedright\arraybackslash}p{(\linewidth - 4\tabcolsep) * \real{0.1364}}
  >{\raggedright\arraybackslash}p{(\linewidth - 4\tabcolsep) * \real{0.6136}}@{}}
\toprule\noalign{}
\begin{minipage}[b]{\linewidth}\raggedright
Invariant
\end{minipage} & \begin{minipage}[b]{\linewidth}\raggedright
Type
\end{minipage} & \begin{minipage}[b]{\linewidth}\raggedright
Recovers in manifold case
\end{minipage} \\
\midrule\noalign{}
\endhead
\bottomrule\noalign{}
\endlastfoot
Index pairing \(\langle [D], [p]\rangle\) & Integer-valued,
\(\mathbb{Z}\) & Topological invariants (Â-genus, Chern numbers) \\
Spectral distance \(d_D(\phi, \psi)\) & Metric on
\(\mathcal{S}(\mathcal{A})\) & Riemannian geodesic distance \\
\end{longtable}
}

Together they realize Connes' programme: every Riemannian (spin)
manifold is reconstructible from its spectral triple, with both the
topology (via \(K\)-homology) and the metric (via \(d_D\)) encoded in
the spectral data.

\textbf{Formalization note.} Unlike the index pairing --- which lives
naturally in the bounded \(F\)-picture --- the spectral distance is
intrinsically tied to the \emph{unbounded} \(D\). In the bounded picture
the commutator \([F, \pi(a)]\) is \emph{compact} (not bounded with
controlled norm), so \(\|[F, \pi(a)]\|\) does \textbf{not} define the
same Lipschitz seminorm; the metric built from it generally differs from
\(d_D\). The needed unbounded commutator API is now implemented on
\texttt{IsOddSpectralTriple}: \texttt{domainRepresentation} and
\texttt{commutatorOnDomain} define the domain map,
\texttt{exists\_commutator\_bound} proves a global linear norm bound,
and \texttt{boundedCommutator} is the canonical continuous extension,
with \texttt{boundedCommutator\_apply} proving agreement on
\(\mathop{\mathrm{dom}}D\). No parallel structure or extra extension
axiom is needed.

The definition itself is one line once the data is available:
\(d_D(\phi, \psi) \;=\; \sup\{|\phi(a) - \psi(a)| : a^* = a,\quad
\|\texttt{hT.boundedCommutator a}\| \leq 1\}\).

\subsection{2.8 Cyclic cohomology and the Chern--Connes
character}\label{cyclic-cohomology-and-the-chernconnes-character}

A spectral triple represents a \emph{K-homology class} of
\(\mathcal{A}\); the natural algebraic dual to K-homology is
\textbf{cyclic cohomology} \(HC^*(\mathcal{A})\) (Connes 1985). The
image of the triple under this duality is the \textbf{Chern--Connes
character} \[
\mathrm{ch}(D) \;\in\; HP^*(\mathcal{A})
\] in periodic cyclic cohomology. The Phase 2 index pairing factors
through it: \[
\bigl\langle [D],\,[p]\bigr\rangle
\;=\; \bigl\langle \mathrm{ch}(D),\,\mathrm{ch}(p)\bigr\rangle,
\] where \(\mathrm{ch}(p) \in HC_*(\mathcal{A})\) is the Chern character
of the projection on the K-theory side. Cyclic cohomology is precisely
the cohomological substrate that makes the integer pairing of §2.2 into
a Chern--Weil-style integration.

\textbf{Cyclic cohomology in short.} For a unital algebra
\(\mathcal{A}\) over \(\mathbb{C}\), the \textbf{Hochschild complex}
\(C^n(\mathcal{A}, \mathcal{A}^*) = \mathrm{Hom}(\mathcal{A}^{\otimes(n+1)}, \mathbb{C})\)
with coboundary \(b\) computes Hochschild cohomology
\(HH^*(\mathcal{A})\). Restricting to cyclically-symmetric cochains and
adjoining Connes' second differential \(B\) yields the
\((b, B)\)-bicomplex; its total cohomology is \textbf{cyclic cohomology}
\(HC^*(\mathcal{A})\). The periodicity operator \(S\) stabilizes the
degree mod 2 and produces \(\mathbb{Z}/2\)-graded \textbf{periodic
cyclic cohomology} \(HP^*(\mathcal{A})\).

\textbf{Two presentations of \(\mathrm{ch}(D)\).}

\begin{itemize}
\tightlist
\item
  \textbf{Connes' algebraic \((b, B)\)-cocycle} (1985; finitely
  summable, purely algebraic). For a \((p+1)\)-summable graded Fredholm
  module \((\mathcal{A}, \mathcal{H}, F, \gamma)\) with \(p\) even, \[
  \tau_p(a_0, \dots, a_p)
  \;:=\; \lambda_p\,\mathrm{Tr}\bigl(\gamma\,F\,[F, \pi(a_0)]\,[F, \pi(a_1)]\,\cdots\,[F, \pi(a_p)]\bigr)
  \] (with \(\lambda_p\) a normalization) is a cyclic \(p\)-cocycle on
  \(\mathcal{A}\), and its class \([\tau_p] \in HC^p(\mathcal{A})\)
  represents \(\mathrm{ch}(D)\). Summability --- \([F, \pi(a)]\) lies in
  the Schatten ideal \(\mathcal{L}^{p+1}(\mathcal{H})\) for every \(a\)
  --- replaces the dimension axiom of a smooth manifold. \emph{Lives
  on the classical bounded Fredholm-module layer; the corresponding
  Lean structure is planned, and no unbounded \(D\) or heat semigroup
  is needed once that layer exists.}
\item
  \textbf{JLO cocycle} (Jaffe--Lesniewski--Osterwalder 1988; entire,
  unbounded). For an unbounded \(D\), \[
  \varphi_{2n}(a_0, \dots, a_{2n}) \;=\;
  \int_{\Delta_{2n}} \mathrm{Tr}\!\Bigl(\gamma\,\pi(a_0)\,e^{-s_0 D^2}\,[D, \pi(a_1)]\,e^{-s_1 D^2}\,\cdots\,[D, \pi(a_{2n})]\,e^{-s_{2n} D^2}\Bigr)\,ds
  \] on the standard \(2n\)-simplex \(\Delta_{2n}\). The family
  \((\varphi_{2n})_n\) assembles to an \textbf{entire cyclic cocycle}
  representing \(\mathrm{ch}(D)\). \emph{Needs the heat semigroup
  \(e^{-tD^2}\) and hence the unbounded \(D\); trades finite summability
  for entire-cyclic regularity.}
\end{itemize}

Both presentations represent the same class in \(HP^*(\mathcal{A})\)
when both are defined, and both deliver the same integer pairing with
\(K_0(\mathcal{A})\).

\textbf{Manifold specialization --- HKR / Connes' theorem.} For
\(\mathcal{A}= C^\infty(M)\) on a closed smooth manifold, the
\textbf{Hochschild--Kostant--Rosenberg theorem} (continuous version)
gives \(HH^n(C^\infty(M)) \cong \Omega^n(M)\), and Connes' refinement
promotes this to the periodic cyclic statement \[
HP^*\bigl(C^\infty(M)\bigr) \;\cong\; H^*_{\mathrm{dR}}(M;\, \mathbb{C}),
\qquad \text{$\mathbb{Z}/2$-graded by parity of degree.}
\] Under this iso, \(\mathrm{ch}(D_M)\) becomes the index density
\(\widehat{A}(M) \in H^*_{\mathrm{dR}}(M;\mathbb{C})\), and pairing with
\(\mathrm{ch}(p_E)\) for a vector bundle \(E\) recovers the
Atiyah--Singer integrand \(\widehat{A}(M) \cup \mathrm{ch}(E)\). The
square \[
\begin{array}{ccc}
K_0(\mathcal{A}) \otimes K^0(\mathcal{A}) & \xrightarrow{\langle\cdot,\cdot\rangle} & \mathbb{Z}\\
\mathrm{ch}\otimes\mathrm{ch}\;\downarrow & & \downarrow \\
HP_*(\mathcal{A}) \otimes HP^*(\mathcal{A}) & \xrightarrow{\langle\cdot,\cdot\rangle} & \mathbb{C}
\end{array}
\] commutes: the integer index lifts to a complex pairing that happens
to land in \(\mathbb{Z}\) for spectral-triple data.

\textbf{Direct alternative on manifolds.} Cohomology of \(M\) can also
be read off the \emph{specific} Hodge--de Rham triple of §2.6: with
\(\mathcal{H}= L^2(\Lambda^* M)\), \(D = d + d^*\),
\(\gamma = (-1)^{\deg}\), Hodge theory gives
\(\ker(D)|_{\Lambda^k} \cong H^k_{\mathrm{dR}}(M;\,\mathbb{R})\). This
is concrete but tied to one specific triple, and is \emph{not} the
NCG-native answer: it computes cohomology as a kernel of an operator on
forms rather than as a structure intrinsic to \(\mathcal{A}\). The
cyclic-cohomology picture is the framework-level answer, computing
manifold cohomology as a derived invariant of \(\mathcal{A}\) alone.

\textbf{Formalization route.} The algebraic \((b, B)\)-cocycle is the
natural target: it lives on the planned bounded Fredholm-module layer
and requires no unbounded-operator API, only Schatten-class membership of
the commutators. Cyclic cohomology of \(\mathcal{A}\) itself is a purely
algebraic Hochschild-style construction. The JLO presentation and the
HKR / Connes iso to de Rham cohomology are heavier and defer to the
unbounded-API phase / Phase 4a respectively. See Phase 2.75 in §3.6 for
the encoding plan.

\subsection{2.9 Genuinely noncommutative
instances}\label{genuinely-noncommutative-instances}

The framework's axioms (§2.1) are stated for an arbitrary unital
\(*\)-algebra \(\mathcal{A}\); they do not require commutativity.
Several classes of genuinely noncommutative examples fit the implemented
unbounded core and should fit the planned bounded Fredholm-module layer
without changing the algebra typeclasses of §3.2.

\textbf{Matrix algebras and finite triples.}
\(\mathcal{A}= M_n(\mathbb{C})\) acting on
\(\mathcal{H}= \mathbb{C}^n \otimes \mathbb{C}^n\) (or its
\(\mathbb{Z}/2\)-graded version) gives the simplest finite-dimensional
even spectral triple. Two-point spaces and finite-dimensional models for
gauge theories (Connes--Lott; Chamseddine--Connes spectral standard
model) live in this class. These are also the cleanest examples for
testing the cyclic-cohomology Chern character of §2.8 in a setting where
all sums are finite.

\textbf{The noncommutative torus \(A_\theta\).} For
\(\theta \in [0, 1) \setminus \mathbb{Q}\), the \textbf{irrational
rotation algebra} \(A_\theta\) is the universal \(C^*\)-algebra
generated by unitaries \(U, V\) with \(UV = e^{2\pi i \theta}\,VU\); its
smooth subalgebra \(A_\theta^\infty\) (rapid-decay Fourier coefficients)
is the pre-\(C^*\) algebra for the spectral triple. Connes (1980)
constructs a canonical even spectral triple
\((A_\theta^\infty,\, \mathcal{H}= \ell^2(\mathbb{Z}^2) \otimes \mathbb{C}^2,\, D, \gamma)\)
modelled on the flat-torus Dirac operator: with \(\tau(a) := a_{0,0}\)
the unique trace, \(\mathcal{H}\) is the GNS Hilbert space tensored with
the spinor representation, and
\(D = \delta_1 \otimes \sigma_1 + \delta_2 \otimes \sigma_2\) where
\(\delta_1, \delta_2\) are the canonical derivations dual to \(U, V\).
Strikingly \(D\) has the \emph{same spectrum} as the flat \(T^2\) Dirac
(Connes' isospectral deformation) --- the geometry sees the deformation
only through the algebra, not through the eigenvalues.

The K-theory is \(K_0(A_\theta) \cong \mathbb{Z}^2\) as an
abstract abelian group, while the \textbf{range of the trace} on
projections is
\(\tau\bigl(K_0(A_\theta)\bigr) = \mathbb{Z}+ \theta\mathbb{Z}\subset \mathbb{R}\)
(Pimsner--Voiculescu). The \textbf{Powers--Rieffel projection}
\(p_\theta \in A_\theta\) realizes \(\tau(p_\theta) = \theta\), so the
trace pairing
\(\langle \tau, [p_\theta]\rangle = \theta \in \mathbb{R}\setminus \mathbb{Q}\)
is \textbf{irrational}. This is a different pairing from the Fredholm
pairing:

\begin{itemize}
\tightlist
\item
  pairing a K-homology/Fredholm class with \([p]\in K_0(A_\theta)\)
  always gives an integer; for the standard fundamental class and the
  Rieffel generator it gives the integral Chern number, up to orientation;
\item
  pairing the degree-0 cyclic cocycle \(\tau\) with the same K-theory
  class gives the real number \(\tau(p)\), in particular \(\theta\) for
  \(p_\theta\).
\end{itemize}

Cyclic cohomology accommodates both the trace cocycle and the higher
cocycle representing the fundamental class, but it does not turn a
Fredholm index into an irrational number. The two pairings are
complementary and must remain distinct.

\textbf{Group \(C^*\)-algebras.} For a finitely generated discrete group
\(\Gamma\) with a proper length function
\(\ell\colon \Gamma \to \mathbb{R}_{\geq 0}\) (word length on a Cayley
graph; restriction of a Lie-group norm), the reduced \(C^*\)-algebra
\(C^*_r(\Gamma)\) has a spectral triple with
\(\mathcal{H}= \ell^2(\Gamma)\), \(\pi\) left convolution, and
\(D\colon \delta_g \mapsto \ell(g)\,\delta_g\) the length operator.
Compact resolvent holds iff \(\ell\) is \textbf{proper}:
\(\#\{g : \ell(g) \leq R\} < \infty\). This is the basic input to the
Baum--Connes assembly map and to Connes--Moscovici's transverse
fundamental class for foliations.

\textbf{Quantum groups.} \(C^*\)-completions of
\(\mathcal{O}(SU_q(2))\), \(\mathcal{O}(SU_q(n))\), and Podleś quantum
spheres carry spectral triples (Chakraborty--Pal; Connes--Landi;
Dąbrowski--Sitarz). These deform the canonical Dirac triples of compact
Lie groups / homogeneous spaces, and exhibit the cyclic-cohomology
refinement explicitly via twisted traces.

\textbf{Crossed products \(\mathcal{A}\rtimes G\).} Classically,
suitable equivariant Fredholm data can be reorganized as non-equivariant
data over a crossed product, subject to the usual covariance and
completion hypotheses. This is the algebraic mechanism by which
equivariance is absorbed into a larger algebra.

\textbf{Fuzzy spheres.} \(\mathcal{A}= M_n(\mathbb{C})\) acting on a
spin-\(j\) representation of \(\mathfrak{su}(2)\) with \(j = (n-1)/2\),
equipped with a \(D\) built from \(\mathfrak{su}(2)\) Casimirs, gives a
sequence of finite-dimensional spectral triples converging --- in
Rieffel's quantum Gromov--Hausdorff sense (§6.6) --- to the spectral
triple of \(S^2\) as \(n \to \infty\). The canonical \emph{finite} test
case for the continuum-limit programme; closer to §6 stretch goals than
to initial deliverables.

\textbf{Implications for the formalization.} None of these examples
requires changes to \texttt{Basic.lean}: its assumptions are
\texttt{Semiring\ A}, \texttt{StarRing\ A}, and
\texttt{Algebra\ K\ A}; no \texttt{StarModule\ K\ A} is required. The
noncommutative torus is
concrete enough --- purely algebraic on \(\ell^2(\mathbb{Z}^2)\), no
spin bundles, no Sobolev embedding --- to formalize at Phase 3 level;
see Phase 3.6 in §3.6. Group \(C^*\)-algebras, quantum groups, and fuzzy
spheres are stretch goals comparable in scope to the manifold
construction.

\section{3. Formalization design}\label{formalization-design}

\subsection{3.1 Encoding choice: bounded vs unbounded
D}\label{encoding-choice-bounded-vs-unbounded-d}

The unbounded \texttt{LinearPMap} formulation is the implemented spine.
The following layers are complete: the odd/even spectral-triple core in
\texttt{Basic.lean}; resolvent definitions in \texttt{Resolvent.lean};
the off-real-axis self-adjoint resolvent criterion and
\texttt{IsCompactResolventSpectralTriple} in
\texttt{CompactResolvent.lean}; the bounded predicate
\texttt{Fredholm.IsFredholm} and \texttt{Fredholm.index}; the compact
perturbation theorem in \texttt{CompactOperators.lean}; and the
full-operator \texttt{gradedKernelIndex} in \texttt{Index.lean}.

The chiral restriction and projection pairing are not yet defined.
Finite or \(p\)-summability is also open: compact resolvent is strictly
weaker than Schatten-class summability. \texttt{FinitelySummable.lean}
is only a compatibility import and
\texttt{IsFinitelySummableSpectralTriple} is deprecated.

The remaining large bridge is the bounded transform
\(F=D(1+D^2)^{-1/2}\) for unbounded self-adjoint \(D\). Classically it
preserves the index pairing, but that functional-calculus theorem is not
yet formalized. Spectral distance remains on the unbounded spine.

\subsection{3.2 Algebra category}\label{algebra-category}

We parametrize over \texttt{(K\ :\ Type*)\ {[}RCLike\ K{]}} (so the
constructions work over both \(\mathbb{R}\) and \(\mathbb{C}\)) and an
algebra
\texttt{(A\ :\ Type*)\ {[}Semiring\ A{]}\ {[}StarRing\ A{]}\ {[}Algebra\ K\ A{]}}.
The implemented core does not require a \texttt{StarModule\ K\ A}
instance.

We deliberately do \textbf{not} require \texttt{A} to be a
\texttt{C*}-algebra at this level: the representation field
\texttt{π\ :\ A\ →⋆ₐ{[}K{]}\ (H\ →L{[}K{]}\ H)} records a
\(*\)-homomorphism into \(\mathcal{B}(\mathcal{H})\), which is the only
place boundedness of \(\pi(a)\) enters. This matches Connes'
``pre-\(C^*\)'' stance: \(\mathcal{A}\) may be the dense subalgebra
\(C^\infty(M)\) of \(C(M)\).

\subsection{3.3 Hilbert space}\label{hilbert-space}

\texttt{(H\ :\ Type*)\ {[}NormedAddCommGroup\ H{]}\ {[}InnerProductSpace\ K\ H{]}\ {[}CompleteSpace\ H{]}}.

Completeness is essential for compactness arguments. Separability is not
imposed at the type level (it is rarely needed in Mathlib formalizations
of operator theory and can be added when required).

\subsection{3.4 Grading}\label{grading}

The grading \(\gamma\) is bounded, so it lives in
\texttt{H\ →L{[}K{]}\ H}; the unbounded Dirac operator is a separate
parameter. \texttt{IsEvenSpectralTriple\ A\ D\ π\ γ} extends the odd
core with self-adjointness, \(\gamma^2=1\), commutation with the
representation, preservation of \texttt{D.domain}, and domain-sensitive
anticommutation. The code proves the \(\pm1\) eigenspace decomposition
and preservation of \texttt{Dkernel}. A future bounded even Fredholm
module may mirror this design, but no such structure exists today.

\subsection{3.5 File structure}\label{file-structure}

\begin{verbatim}
SpectralTriples.lean                       -- root re-export
SpectralTriples/
├── Basic.lean                             -- unbounded odd/even core and bounded commutators
├── Resolvent.lean                         -- LinearPMap resolvent definitions
├── CompactResolvent.lean                  -- criterion and compact-resolvent package
├── FinitelySummable.lean                  -- deprecated compatibility import
├── Fredholm.lean                          -- bounded IsFredholm and integer index
├── CompactOperators.lean                  -- compact-adjoint and 1 − K Fredholm theory
├── Index.lean                             -- Dkernel and gradedKernelIndex
├── DiagonalOperator.lean                  -- block-diagonal ℓ² operators
├── FourierHolomorphic.lean                -- automorphic function-space dimension results
├── HermiteL2.lean                         -- normalized Hermite L² orthonormal family
└── Examples/
    ├── Circle.lean                        -- S¹ Fourier spectral triple
    ├── Shift.lean                         -- bounded unilateral shift index
    ├── MagneticDirac.lean                 -- bounded magnetic phase and translations
    ├── ThetaSections.lean                 -- explicit automorphic theta sections
    └── Torus.lean                         -- flat T² Fourier spectral triple
\end{verbatim}

Planned modules such as \texttt{BoundedTransform.lean},
\texttt{Distance.lean}, \texttt{Cyclic.lean},
\texttt{Examples/TorusLineBundle.lean}, and the \texttt{Manifold/}
hierarchy are not part of the checked-in tree.

\subsection{3.6 Phases}\label{phases}

\textbf{Phase 1 --- abstract definitions and compact resolvent. [DONE]}
\texttt{Basic.lean} implements the unbounded odd/even core and the
canonical bounded commutator extension. \texttt{Resolvent.lean} supplies
the resolvent definitions. \texttt{CompactResolvent.lean} proves the
formerly open basic criterion
\[
D=D^*,\quad \operatorname{Im}z\ne0 \Longrightarrow z\in\rho(D),
\]
including the closed- and dense-range steps, and packages
\texttt{IsCompactResolventSpectralTriple}. Circle and torus instantiate
it. This is compact resolvent, not \(p\)-summability.

\textbf{Phase 2 --- Fredholm foundations and index pairing. [PARTIAL]}
The implemented foundation consists of \texttt{Fredholm.IsFredholm},
\texttt{Fredholm.index}, the \(1-K\) compact-perturbation theorem,
finite-dimensionality of the full Dirac kernel under compact resolvent,
and \texttt{gradedKernelIndex}. The unilateral shift and bounded
magnetic phase have explicit Fredholm and index theorems.

\texttt{gradedKernelIndex} is not yet the index of a formalized chiral
operator. Remaining work is to define the graded spaces and
domain-restricted \(D^+\), prove it Fredholm, define \(pD^+p\), and prove
agreement with the graded-kernel formula. Direct sums, homotopy
invariance, and K-theory descent are downstream. No fill-in partial
isometry may be assumed without proving its defect dimensions and
Fredholm properties.

\textbf{Phase 2.5 --- Connes' spectral distance (see §2.7).
[FOUNDATION DONE; DISTANCE PLANNED]} The unbounded core now supplies the
canonical continuous extension \texttt{boundedCommutator}; the former
commutator-extension blocker is closed.

\begin{itemize}
\tightlist
\item
  Define \texttt{connesDistance\ (φ\ ψ\ :\ StateSpace\ A)\ :\ ℝ≥0∞} as
  the supremum of \(|\phi(a) - \psi(a)|\) over self-adjoint
  \(a \in \mathcal{A}\) with
  \(\|\texttt{hT.boundedCommutator a}\| \leq 1\).
\item
  Prove: the Connes distance on a commutative finite spectral triple (a
  graph, via the §3.9 \texttt{SusyGraphop} bridge) equals the
  shortest-path graph distance. This is the simplest non-trivial
  verification and re-validates the §3.9 bridge.
\item
  (Goal, beyond initial scope.) Connes' theorem on the canonical Dirac
  triple: \(d_D(\delta_x, \delta_y) = d_g(x, y)\) --- geodesic distance
  recovery. Requires the manifold construction (Phase 4) and significant
  analytic work; flag as Phase 4 territory.
\end{itemize}

\textbf{Phase 2.75 --- cyclic cohomology and the Chern--Connes character
(see §2.8).} The cohomological dual to Phase 2: encode the cocycle on
the algebra side that the index pairing factors through.

\begin{itemize}
\tightlist
\item
  \texttt{Cyclic.lean}: build Hochschild cohomology
  \(HH^*(\mathcal{A})\) and cyclic cohomology \(HC^*(\mathcal{A})\) from
  the algebraic \((b, B)\)-bicomplex of an arbitrary unital
  \(*\)-algebra; define the periodicity operator \(S\) and the
  \(\mathbb{Z}/2\)-graded periodic cyclic cohomology
  \(HP^*(\mathcal{A})\). Purely algebraic --- Mathlib's cochain-complex
  and tensor-power infrastructure suffices.
\item
  \texttt{ChernCharacter.lean}: on a \emph{finitely-summable} graded
  \texttt{FredholmModule} (parametrized by \(p\) via an added
  Schatten-class axiom on \([F, \pi(a)]\)), define the Connes
  \((b, B)\)-cocycle \(\tau_p\) and prove it is a cyclic \(p\)-cocycle
  on \(\mathcal{A}\). Define the K-theory Chern character
  \(\mathrm{ch}(p) \in HC_*(\mathcal{A})\) for a projection
  \(p \in \mathcal{A}\). State and prove
  \[\bigl\langle [F],\,[p] \bigr\rangle \;=\; \bigl\langle \mathrm{ch}(F),\,\mathrm{ch}(p) \bigr\rangle,\]
  the factorization of the integer index pairing through cyclic
  cohomology.
\item
  \textbf{Schatten ideals are partial in Mathlib.} A
  \texttt{pSummableFredholmModule\ p} extension carrying the axiom
  \([F, \pi(a)] \in \mathcal{L}^{p+1}(\mathcal{H})\) for every \(a\) is
  the cleanest design, but Schatten classes themselves are still being
  upstreamed (\texttt{Mathlib.Analysis.InnerProductSpace.Schatten}).
  Design Phase 2.75 to be parametric in that gap: require Schatten
  membership as an explicit theorem parameter or a centrally declared,
  policy-vetted interface. Local \texttt{sorry}/\texttt{admit}
  placeholders are prohibited.
\item
  \textbf{Deferred:} the JLO presentation (needs unbounded \(D\) and the
  heat semigroup; out of scope until the unbounded-API phase) and the
  HKR / Connes iso \(HP^*(C^\infty(M)) \cong H^*_{\mathrm{dR}}(M)\)
  (needs smooth-form infrastructure; out of scope until Phase 4a / 5).
\end{itemize}

This phase is \textbf{algebraic-only}: zero unbounded analysis, zero
manifold tools. It states the framework's central structural theorem ---
\emph{the index pairing factors through the Chern--Connes character} ---
at the level of generality of Phase 2, lifting the integer pairing to
its natural cohomological refinement.

\textbf{Phase 3 --- example: the circle \(S^1\). [CORE DONE]}
\texttt{Examples/Circle.lean} builds the Dirac triple
on \(S^1\) directly from Fourier series. With
\(\mathcal{H}= \ell^2(\mathbb{Z})\) (Fourier basis of \(L^2(S^1)\)) and
\(D= -i\,d/d\theta\) acting as the diagonal operator
\(\delta_n \mapsto n\,\delta_n\), the spectrum is \(\mathbb{Z}\), the implemented
resolvent \((i\,1-D)^{-1}\) is the diagonal operator with eigenvalues
\(1/(i-n)\) (compact because \(|1/(i-n)| \to 0\)). The implemented algebra is the
star-subalgebra generated by the Fourier shift \texttt{W 1}, the operator image of the
trigonometric/Laurent polynomials; bounded commutators are proved on the generators and
propagated algebraically. The full classical \(C^\infty(S^1)\) algebra is not bundled here.
This bypasses spin bundles, Sobolev
compactness, and Levi--Civita connections. The circle is
also the natural testbed for Phase 4.5 (the rotation action of \(S^1\)
on itself). Caveat: \(S^1\) is odd-dimensional, so the Dirac triple is
\emph{ungraded} (no chirality \(\gamma\)) --- it pairs with
\(K^1(\mathcal{A})\), not \(K_0(\mathcal{A})\). The index pairing here
is the Toeplitz-type spectral-flow pairing rather than the even pairing
of Phase 2. \texttt{Examples/Shift.lean} separately proves the bounded
unilateral forward shift Fredholm with index \(-1\); that theorem has
not yet been assembled into a \(K^1\) pairing for the circle triple.

\textbf{Phase 3.5 --- example: \(T^2\) with a degree-\(d\) line bundle.
[PARTIAL]}
The flat torus \(T^2 = \mathbb{R}^2/\mathbb{Z}^2\) with the trivial spin
structure is the simplest \emph{even-dimensional} concrete Dirac triple,
and the simplest case where the index pairing produces a non-trivial,
parameter-controlled integer.

\emph{Setup.} The spinor bundle on \(T^2\) is rank 2 with both \(S^\pm\)
trivial (since \(TT^2\) is trivial). In the standard complex coordinate
\(z = x_1 + i x_2\) the Dirac operator reduces to
\(D^+ \cong 2\,\overline{\partial}\colon C^\infty(T^2) \to \Omega^{0,1}(T^2)\).
At a nonzero Fourier mode \(e^{2\pi i(m x_1 + n x_2)}\), the
\(2\times2\) Dirac block has the two simple eigenvalues
\(\pm2\pi\sqrt{m^2 + n^2}\).
\(\mathcal{H}= L^2(T^2, S) \cong \ell^2(\mathbb{Z}^2; \mathbb{C}^2)\)
carries these blocks diagonally.

\emph{Untwisted index.}
\(\mathrm{ind}(D^+) = h^0(\mathcal{O}_{T^2}) - h^1(\mathcal{O}_{T^2})
= 1 - 1 = 0\) (genus-1 Hodge). Baseline.

\emph{Twist by a line bundle \(L\) of degree \(d\).} On \(T^2\), line
bundles are classified by
\(\mathrm{Pic}(T^2) \to H^2(T^2; \mathbb{Z}) \cong \mathbb{Z}\), with
the degree map \(L \mapsto c_1(L)[T^2] = d\). The twisted Dirac
\(D_L^+ \colon \Gamma(S^+ \otimes L) \to \Gamma(S^- \otimes L)\)
satisfies Atiyah--Singer / Riemann--Roch: \[
\mathrm{ind}(D_L^+) \;=\; \chi(L) \;=\; \deg L - g + 1 \;=\; d.
\] This is the cleanest sanity check available for the index pairing: a
single integer, parametrized by \(\mathbb{Z}\), computable in several
independent ways. The next paragraph spells out the framework-level
computation.

\emph{Classical worked computation: the index is \(d\), three ways.}
Only the untwisted Fourier triple---using the star-subalgebra generated by the two
coordinate shifts, not the full \(C^\infty(T^2)\) algebra---and the supporting
function/phase results identified below are currently formalized.

\textbf{Base spectral data.} With \(\mathcal{A}= C^\infty(T^2)\),
\(\mathcal{H}= L^2(T^2) \otimes \Sigma\) (\(\Sigma = \mathbb{C}^2\)),
\(D = -i(\sigma_1\,\partial_1 + \sigma_2\,\partial_2)\),
\(\gamma = \sigma_3\). Fourier diagonalises:
\(\mathcal{H}\cong \ell^2(\mathbb{Z}^2) \otimes \Sigma\), with
\(\pi(a)\) acting by convolution by the Fourier coefficients of \(a\)
and \(D\) diagonal on Fourier modes, \[
D|_{(m, n)} \;=\; 2\pi \begin{pmatrix} 0 & m - in \\ m + in & 0 \end{pmatrix}
\quad \text{on } \Sigma.
\] Classically, the bounded transform \(F = D(1 + D^2)^{-1/2}\) has off-diagonal
block \(F^+_{(m, n)} = 2\pi(m + in)/\sqrt{1 + 4\pi^2(m^2 + n^2)}\),
vanishing only at \((m, n) = (0, 0)\). The untwisted check is immediate:
\(\ker F^+ = \ker F^- = \mathbb{C}\), so
\(\langle [F], [1]\rangle = 1 - 1 = 0\).

\textbf{Serre--Swan projection \(p_d\) (TODO: construct from clutching
data).} Present the degree-\(d\) line bundle as the quotient of
\([0,1]\times S^1\times\mathbb{C}\) by
\[
(1,y,z)\sim(0,y,e^{2\pi i d y}z).
\]
Choose a finite open cover \(\{U_i\}\), unitary transition functions
\(g_{ij}\), and a subordinate smooth partition of unity \(\{\rho_i\}\).
The standard construction gives
\[
(p_d)_{ij}(x)=\sqrt{\rho_i(x)}\,g_{ij}(x)\,\sqrt{\rho_j(x)}.
\]
The cocycle identity and \(\sum_i\rho_i=1\) prove \(p_d^2=p_d\), and
unitarity proves \(p_d^*=p_d\). The Lean TODO is to choose the cover and
partition explicitly and prove, with the orientation/sign convention
stated, that the image bundle has Chern number \(d\). Do not reuse the earlier sine/cosine \(2\times2\) formula: its
image admits a periodic nowhere-zero global frame and therefore has
degree zero.

\textbf{Compressed operator (not yet in Lean).} Amplify the Fredholm module to act on
\(\mathcal{H}' := \mathcal{H}\otimes \mathbb{C}^N_{\mathrm{aux}}\) via
\(\pi_N := \pi \otimes \mathrm{id}_{M_N}\), \(F' := F \otimes 1_N\),
\(\gamma' := \gamma \otimes 1_N\). Let \(P_d := \pi_N(p_d)\) --- a bounded
self-adjoint idempotent on \(\mathcal{H}'\). The pairing is the Fredholm
index of the compression \[
T_d^+ \;:=\; P_d\,(F^+ \otimes 1_N)\,P_d \;\colon\; P_d\,\mathcal{H}^{\prime,+} \to P_d\,\mathcal{H}^{\prime,-}.
\] Serre--Swan identifies
\(P_d\,\mathcal{H}^{\prime,\pm} \cong L^2(T^2,\,L_d \otimes \Sigma^\pm)\),
and under this iso \(T_d^+\) is the bounded transform of the twisted
Dirac
\(D_{L_d}^+ \colon L^2(T^2, L_d \otimes \Sigma^+) \to L^2(T^2, L_d \otimes \Sigma^-).\)
The bounded transform preserves \(\ker\), \(\mathrm{coker}\), and the
Fredholm index, so \[
\bigl\langle [F],\,[p_d] \bigr\rangle \;=\; \mathrm{ind}(T_d^+) \;=\; \mathrm{ind}(D_{L_d}^+).
\]

\textbf{Three independent computations of \(d\).}

\begin{itemize}
\tightlist
\item
  \textbf{(i) Riemann--Roch.}
  \(\mathrm{ind}(D_{L_d}^+) = \chi(L_d) = \deg L_d + 1 - g(T^2) = d + 1 - 1 = d\).
\item
  \textbf{(ii) Aharonov--Casher / Landau levels.} Choose the
  constant-curvature connection on \(L_d\), with curvature
  \(F_{L_d} = 2\pi i d\,dx_1 \wedge dx_2\) (total flux \(2\pi d\)). The
  Lichnerowicz--Bochner identity on the flat \(T^2\) specialises to
  \(D_{L_d}^2 = \nabla^*\nabla - \pi d\,\sigma_3.\) For \(d > 0\): on
  \(\Sigma^+\), \(\ker\) requires \(\nabla^*\nabla = \pi d\) --- the
  lowest Landau level, with degeneracy exactly \(d\); on \(\Sigma^-\),
  \(\ker\) would require \(\nabla^*\nabla = -\pi d < 0\), impossible.
  Thus \(\dim\ker D_{L_d}^+ = d\) and \(\dim\ker D_{L_d}^- = 0\), giving
  \(\mathrm{ind} = d\). The chiralities swap for \(d < 0\) and the index
  is still \(d\).
\item
  \textbf{(iii) Chern--Connes character (Phase 2.75).} The degree-2
  cyclic cocycle of the \(T^2\) Fredholm module pairs with the K-theory
  Chern character of \(p_d\) as \[
  \bigl\langle [F],\,[p_d] \bigr\rangle \;=\; \bigl\langle \mathrm{ch}_2(F),\,\mathrm{ch}_2(p_d) \bigr\rangle
  \;=\; \frac{1}{2\pi i}\int_{T^2} \mathrm{tr}\bigl(p_d\,dp_d \wedge dp_d\bigr)
  \;=\; c_1(L_d)[T^2] \;=\; d.
  \] The middle integral is the classical Chern--Weil formula for the
  first Chern class of the rank-1 image of a matrix-valued projection.
  The pairing factors literally through the de Rham class
  \(c_1 \in H^2(T^2;\,\mathbb{R})\) --- Phase 2.75 in action on a
  manifold.
\end{itemize}

Three perspectives, same integer. \textbf{(ii)} is the direct
operator-theoretic verification on \(\ell^2(\mathbb{Z}^2)\) and is the
most natural target for \texttt{Examples/TorusLineBundle.lean};
\textbf{(iii)} is the framework-native one and is the natural target for
the Phase 2.75 deliverable. Both pass through the bounded-transform
identification \(T_d^+ \leftrightarrow D_{L_d}^+\).

\emph{Status and deliverables.}

\begin{itemize}
\tightlist
\item
  \textbf{Done:} \texttt{Examples/Torus.lean} builds the untwisted
  unbounded even spectral triple, proves compact resolvent, and proves
  \texttt{isEvenSpectralTriple.gradedKernelIndex = 0}. This is not a
  compressed Fredholm pairing theorem.
\item
  \textbf{Done, function-theoretic:} \texttt{ThetaSections.lean} and
  \texttt{FourierHolomorphic.lean} give the exact positive-degree
  automorphic function dimension and negative-degree vanishing. They do
  not identify operator kernels or cokernels.
\item
  \textbf{Done, bounded model:} \texttt{MagneticDirac.lean} proves
  \texttt{magneticPhase k} Fredholm of index \(k\).
  \texttt{HermiteL2.lean} proves the normalized functions
  \(h_n(x)=H_n(x)e^{-x^2/4}/\sqrt{n!\sqrt{2\pi}}\) form an
  orthonormal family; completeness remains open.
\item
  \textbf{Open:} construct the clutched line bundle, weighted
  \(L^2\)-space, and twisted chiral operator. Identify it with the
  unbounded \(\sqrt{n+1}\)-weighted lowering operator; only that
  operator's bounded polar phase should be identified with
  \texttt{magneticPhase k}.
\end{itemize}

This phase is the natural test bed for Phase 4.5 too (the
\(T^2 = S^1 \times S^1\) action gives an equivariant version with both
index and equivariant index computable).

\textbf{Phase 3.6 --- example: the noncommutative torus \(A_\theta\)
(see §2.9). [PLANNED]} This test case should exhibit two distinct
pairings: an integer Fredholm/Chern-number pairing and the real-valued
trace pairing.

\emph{Setup.} Fix \(\theta \in [0, 1) \setminus \mathbb{Q}\). Encode
\(A_\theta^\infty\) concretely as the algebra of rapid-decay sequences
\(a = (a_{m,n})_{(m,n) \in \mathbb{Z}^2}\) under the twisted convolution
\((ab)_{p,q} = \sum_{m, n} a_{m, n}\,b_{p-m,\,q-n}\,e^{2\pi i \theta\,n(p-m)}\).
The Hilbert space is
\(\mathcal{H}= \ell^2(\mathbb{Z}^2) \otimes \mathbb{C}^2\), the GNS
representation of the unique tracial state \(\tau(a) = a_{0, 0}\)
tensored with the two-dimensional spinor representation. The derivations
\(\delta_j\) act diagonally on Fourier modes:
\(\delta_1(\delta_{m, n}) = 2\pi i m \delta_{m, n}\),
\(\delta_2(\delta_{m, n}) = 2\pi i n \delta_{m, n}\). The Dirac is
\(D = \delta_1 \otimes \sigma_1 + \delta_2 \otimes \sigma_2\), diagonal
on Fourier modes with eigenvalues \(\pm 2\pi\sqrt{m^2 + n^2}\) ---
identical to the flat \(T^2\) Dirac spectrum (Connes' isospectral
deformation). The chirality is \(\gamma = 1 \otimes \sigma_3\).

\emph{Deliverables.}

\begin{itemize}
\tightlist
\item
  \texttt{Examples/NoncommutativeTorus.lean}: construct
  \(A_\theta^\infty\) on rapid-decay \(\mathbb{Z}^2\)-sequences, the GNS
  representation on \(\mathcal{H}\), the Dirac \(D\), and chirality
  \(\gamma\). First verify the unbounded
  \texttt{IsEvenSpectralTriple} and compact-resolvent structures; the
  bounded transform is a later bridge.
\item
  Construct the \textbf{Powers--Rieffel projection}
  \(p_\theta \in A_\theta\) explicitly:
  \(p_\theta = U^* f(V) + g(V) + f(V) U\) for smooth
  \(f, g \in C^\infty(S^1)\) chosen so that
  \(p_\theta^2 = p_\theta = p_\theta^*\) and \(\tau(p_\theta) = \theta\)
  (Rieffel 1981; explicit closed forms for \(f, g\) in terms of bump
  functions).
\item
  Verify \(\tau(p_\theta) = \theta\). The cyclic-cohomology pairing
  \(\langle [\tau], \mathrm{ch}(p_\theta) \rangle\) then equals
  \(\tau(p_\theta) = \theta \in \mathbb{R}\setminus \mathbb{Q}\) --- the
  irrational trace pairing.
\item
  Construct the fundamental Fredholm/K-homology pairing separately and
  compute its integer Chern number on the Rieffel generator. Do not
  identify that integer with the trace value \(\theta\).
\end{itemize}

\emph{Why it matters.} Fredholm indices remain integer-valued for
noncommutative algebras. The canonical trace instead induces
\(K_0(A_\theta)\to\mathbb{R}\) with range
\(\mathbb{Z}+\theta\mathbb{Z}\), while the fundamental higher cocycle
extracts an integral Chern number. Phase 2.75 should encode both and
keep their codomains and cocycles distinct.

\emph{Mathlib gap.} This avoids spin bundles and Sobolev geometry, but
still requires a rapid-decay algebra, twisted convolution, boundedness
of its representation, and direct compactness or Schatten estimates.
None is currently a repository module. Compact resolvent is not a
substitute for a Schatten estimate.

\textbf{Phase 3.7 --- example: the Morse complex of \(S^2\).
\emph{Extra, not core.}} A concrete instance of the Witten--Morse setup
of §2.6, formalizable without semiclassical analysis. Take
\(f \colon S^2 \to \mathbb{R}\) the height function (two critical
points: north pole \(p_+\) of Morse index 2, south pole \(p_-\) of Morse
index 0). The \textbf{Morse complex} is the two-term chain complex
\(\mathbb{Z}\langle p_-\rangle \xrightarrow{0} \mathbb{Z}\langle p_+\rangle\)
with zero differential (no flow lines connect critical points of equal
Morse index parity, and the indices differ by 2). Wrap it as a finite
\texttt{SusyGraphop} with \(n_+ = 1\) (even = index-0 generator),
\(n_- = 0\) (odd = index-1 generators, none), giving \(Q_+ = 0\) and
Witten index = \(\dim\ker Q_+ - \dim\ker Q_- = 1 - 0 = 1\)\ldots{} wait
--- Morse complex of \(S^2\) has one generator each in degree 0 and 2
(both even); odd degree (index 1) is empty. So the Witten index (= Euler
char) is \(1 + 1 = 2\) when we sum over even degrees correctly.
Concretely, set \(n_+ = 2\) (generators \(\langle p_-, p_+\rangle\) in
even degrees 0 and 2), \(n_- = 0\) (no odd generators), \(Q_+ = 0\).
Then Witten index = \(2 - 0
= 2 = \chi(S^2)\). ✓

\emph{Deliverables.}

\begin{itemize}
\tightlist
\item
  \texttt{Examples/MorseS2.lean}: construct the Morse complex of the
  height function on \(S^2\) as a finite \texttt{SusyGraphop} (via the
  \texttt{graphops-qft} bridge of §3.9).
\item
  Verify \(\langle F, [1] \rangle = \chi(S^2) = 2\).
\item
  Compare against the smooth Hodge--de Rham triple of \(S^2\) at
  \(t = 0\) (which has \(\dim \ker = h^0 + h^2 = 2\) on the even side,
  \(h^1 = 0\) on the odd side, so the index is also \(2\)). The point:
  two representatives of the same K-homology class --- smooth (\(t=0\))
  and finite (\(t = \infty\) limit) --- give the same integer.
\end{itemize}

This is a sanity check both for the Witten--Morse framework description
in §2.6 and for the \texttt{SusyGraphop} bridge in §3.9. It does
\emph{not} require formalizing the \(t \to \infty\) analysis itself ---
only the combinatorial endpoint.

\textbf{Phase 4 --- general manifold construction. [PLANNED]} Two sub-phases,
which should not be split into separate Lake targets but should be
tackled in order.

\emph{Phase 4a --- Mathlib lemmas for the spin construction.} Audit
Mathlib for what we need that may not exist yet:

\begin{itemize}
\tightlist
\item
  Spin group \(\mathrm{Spin}(n)\) as a double cover of
  \(\mathrm{SO}(n)\), with the universal property of lifting orthogonal
  frames.
\item
  Complex spinor representation of \(\mathrm{Cl}_n\).
\item
  Levi--Civita connection on the spinor bundle, lifted from \(TM\) via
  the local-frame formula
  \(\nabla^S_X \psi = X(\psi) + \tfrac{1}{4}\sum_{i<j} g(\nabla_X e_i, e_j)\,c(e^i)c(e^j)\,\psi\).
\item
  Sobolev sections \(H^k(M, E)\) for vector bundles on a compact
  manifold, with \textbf{Rellich--Kondrachov compactness} of the
  inclusion \(H^1(M, E) \hookrightarrow L^2(M, E)\).
\end{itemize}

For Rellich--Kondrachov, thread the result as an explicit
\texttt{(rellichKondrachov\ :\ …)} theorem parameter. Any unavoidable
global axiom must be centralized, documented, and policy-vetted. A local
\texttt{sorry} or \texttt{admit} is prohibited. Conditional algebraic
assembly can proceed without obscuring this dependency.

Likely outcome: a separate library of supporting lemmas
(\texttt{Mathlib.Geometry.Manifold.Spin}?) suitable for upstreaming
emerges from this phase.

\emph{Phase 4b --- assembly.} First produce an unbounded
\texttt{IsEvenSpectralTriple} (or odd version) and an
\texttt{IsCompactResolventSpectralTriple} using explicit Phase-4a
hypotheses. Conversion to a bounded Fredholm module and chiral pairing
comes only after the bounded-transform bridge exists.

\textbf{Phase 4.5 --- equivariance. [PLANNED]} The covariance data is
algebraic, but an equivariant chiral index depends on completing Phase 2:
the generic chiral Fredholm operator and its kernel/cokernel are not yet
available.

\emph{Why it's clean.} The unitary \(U(g)\) commutes with \(pF^+p\) by
the covariance axiom, so \(\ker(pF^+p)\) and \(\mathrm{coker}(pF^+p)\)
are \(U(g)\)-invariant. They are therefore finite-dimensional unitary
representations of \(G\), not merely finite-dimensional vector spaces.
The pairing's codomain promotes from \(\mathbb{Z}\) to the
\textbf{representation ring} \(R(G)\) (for compact \(G\)): \[
\mathrm{ind}_G(F) \;:=\; [\ker(pF^+p)] \,-\, [\mathrm{coker}(pF^+p)] \;\in\; R(G).
\] Specializing the character at \(g = 1\) recovers the integer index.

\emph{Fixed-point localization (Atiyah--Segal--Singer 1968).} Evaluating
\(\chi_{\mathrm{ind}_G(F)}(g) = \mathrm{Tr}(g\mid_{\ker}) - \mathrm{Tr}(g\mid_{\mathrm{coker}})\)
gives a global analytic quantity that, by the Atiyah--Segal--Singer
fixed-point theorem, equals a \emph{local} topological sum over the
fixed-point submanifold \(M^g := \{x : g \cdot x = x\}\). For isolated
fixed points this collapses an infinite-dimensional spectral computation
to a finite sum of fractions. (Physics application: the rigorous
foundation of supersymmetric localization --- path integrals collapsing
to finite instanton sums in SUSY field theory.) Formalizing the full
theorem is out of scope, but the \emph{type-level} refinement
\(\mathbb{Z}\rightsquigarrow R(G)\) is in scope and gives the framework
for stating it.

\emph{Crossed-product perspective.} An equivariant spectral triple over
\(\mathcal{A}\) with \(G\)-action \(\sigma\) is mathematically
equivalent to a (non-equivariant) spectral triple over the
\textbf{crossed-product algebra} \(\mathcal{A}\rtimes G\). This is the
NCG mechanism by which an algebra ``knows'' it carries a group action
--- the equivariance is absorbed into a larger algebra. Thus
\texttt{EquivariantFredholmModule\ G\ σ} could equivalently be
reformulated as \texttt{FredholmModule\ ℂ\ (A\ ⋊{[}σ{]}\ G)\ H}; we
choose the mixin form for practical Lean ergonomics, but a
\texttt{crossedProduct\_equivalence} theorem witnessing the two
formulations as the same data is a natural Phase 4.5 deliverable.

\emph{Worked example.} The rotation action \(S^1 \curvearrowright S^1\)
turns the Phase 3 Dirac triple into an \(S^1\)-equivariant one with
explicit Fourier-mode decomposition --- the Fourier modes \(\delta_n\)
are the irreducible \(S^1\)-representations, so the equivariant index
lives in \(R(S^1) = \mathbb{Z}[t, t^{-1}]\) and decomposes by Fourier
weight.

\emph{Mathlib support.} \texttt{Mathlib.RepresentationTheory.*} has
finite group representations, characters, and the representation ring as
a Grothendieck group; the formalization needs no new
representation-theory plumbing beyond what is already upstream.

\emph{Non-compact \(G\) --- Baum--Connes territory.} The representation
ring \(R(G)\) is defined for compact \(G\); for non-compact \(G\) the
right target is \(K_0(C^*_r(G))\), the K-theory of the reduced group
\(C^*\)-algebra, via the Baum--Connes assembly map. This is out of scope
for the initial formalization but flagged here for completeness.

\textbf{Phase 4.75 --- super-Bakry--Émery curvature (see §3.9).
\emph{Extra, not core.}} This phase is a stretch goal: none of the
project's three stated deliverables (abstract triple, Fredholm index
pairing, canonical Dirac triple) requires the super-\(CD(\rho, \infty)\)
machinery, and the Lichnerowicz formula is not load-bearing for the
canonical-triple axioms (those reduce to Chernoff / Rellich--Kondrachov
/ local-frame computations). Pursue Phase 4.75 only if (i) the
\texttt{graphops-qft} unification (§3.9 bridge) is itself a goal, or
(ii) Phase 5 (heat-kernel route to the local index formula) becomes
in-scope. Three concrete deliverables once started:

\begin{itemize}
\tightlist
\item
  \texttt{BakryEmeryFredholmModule\ (ρ\ :\ ℝ)\ extends\ GradedFredholmModule}
  carrying the operator-ordering inequality
  \(\Gamma_2(a, a) \geq \rho\,\Gamma(a, a)\), with \(\Gamma, \Gamma_2\)
  defined via \(\delta(a) = [F, \pi(a)]\) and \(L = F^2\) in the bounded
  picture (or via \(D\) once unbounded API arrives).
\item
  Equivalence theorem: the operator \(CD(\rho, \infty)\) inequality
  \(\Leftrightarrow\) the gradient-decay form \(\|\nabla P_t a\|^2 \leq
  e^{-2\rho t} P_t\|\nabla a\|^2\).
\item
  Bridge \texttt{Bridge/SusyGraphop.lean}: a constructor
  \texttt{ofSusyGraphop\ :\ SusyGraphop\ →\ FredholmModule\ ℝ\ A\ H}
  exhibiting the finite graph case (\texttt{graphops-qft}) as an
  instance, with \(\mathrm{Witten\,index} = \langle [F], [1] \rangle\).
\end{itemize}

The Lichnerowicz specialization (Phase 5 / Phase 4b territory):
\texttt{Ric\_g\ ≥\ K} ⇒ canonical Dirac triple satisfies
\texttt{BakryEmeryFredholmModule\ (K/4)}. Proving this on Mathlib
requires the Lichnerowicz identity for the spinor bundle, which goes
through Phase 4a.

\textbf{Phase 5 --- toward the local index formula.} Out of initial
scope, but worth scoping: Connes--Moscovici residue formula, requiring
zeta functions of \(|D|^{-s}\) and the Wodzicki residue. Connection to
Phase 4.75: the heat-kernel route to Atiyah--Singer (McKean--Singer, via
\(\mathop{\mathrm{Tr}}(\gamma e^{-tD^2})\)) uses the same Lichnerowicz
formula that drives the Phase 4.75 Ricci bound; both are downstream of
the spin-bundle Bochner formula.

\textbf{Phase 6 --- families / Kasparov modules.} Generalize from
Hilbert spaces to Hilbert \(C^*\)-modules over a coefficient
\(C^*\)-algebra \(B\) (see §3.8). Define \texttt{KasparovModule\ A\ B}
in the bounded form, verify it specializes to \texttt{FredholmModule}
when \(B = \mathbb{C}\), and lift the index pairing to a homomorphism
\(K_0(\mathcal{A}) \to K_0(B)\). The Atiyah--Singer index theorem for
families (with \(B = C(X)\)) is the target. Prerequisite: Mathlib
support for Hilbert \(C^*\)-modules and the
\(\mathcal{K}_B(\mathcal{E})\) ideal is sufficient.

\subsection{3.7 Risks and dependencies}\label{risks-and-dependencies}

{\def\LTcaptype{none} % do not increment counter
\begin{longtable}[]{@{}
  >{\raggedright\arraybackslash}p{(\linewidth - 2\tabcolsep) * \real{0.4016}}
  >{\raggedright\arraybackslash}p{(\linewidth - 2\tabcolsep) * \real{0.5984}}@{}}
\toprule\noalign{}
\begin{minipage}[b]{\linewidth}\raggedright
Risk
\end{minipage} & \begin{minipage}[b]{\linewidth}\raggedright
Mitigation
\end{minipage} \\
\midrule\noalign{}
\endhead
\bottomrule\noalign{}
\endlastfoot
Unbounded operator API & Sufficient for the core and compact resolvent;
the bounded-transform functional calculus and general chiral pairing
remain open \\
Spinor bundle / Dirac operator not in Mathlib & Build a self-contained
\texttt{Manifold/} sub-library; upstream candidates \\
Rellich--Kondrachov in the required form & Pass an explicit theorem
hypothesis; any global axiom must be centralized and policy-vetted.
Local \texttt{sorry}/\texttt{admit} placeholders are prohibited \\
Sign / phase conventions diverge between GBF, Connes, Higson--Roe
(e.g.~orientation of \(\gamma\), choice of \(i\) in \(D \mapsto F\)) &
Fix a single convention in \texttt{Basic.lean} docstrings; record
translation tables to the other two references when relevant \\
Hilbert \(C^*\)-module / \(\mathcal{K}_B\) API partial & Defer families
to Phase 6; design Phase 1 to be parametric in \texttt{H} so the lift is
local \\
\end{longtable}
}

\subsection{3.8 Generalizations: encoding equivariance and
families}\label{generalizations-encoding-equivariance-and-families}

The two extensions of §2.4 differ sharply in formalization cost.

\subsubsection{Equivariance --- cheap, add as a
mixin}\label{equivariance-cheap-add-as-a-mixin}

Define a structure parameterized by a topological group \texttt{G}
together with a continuous action on \texttt{A}:

\begin{verbatim}
structure EquivariantFredholmModule
    (G : Type*) [Group G] [TopologicalSpace G] [TopologicalGroup G]
    (σ : G →* (A ≃⋆ₐ[K] A))
    extends FredholmModule K A H where
  U : G →* (H ≃ₗᵢ[K] H)        -- unitary representation
  /-  Strong-operator-topology (SOT) continuity: g ↦ U(g) x is continuous
      for every x ∈ H. Operator-norm continuity is far too strong here —
      for the typical regular representation of an infinite group on
      L²(H) the orbit map is never norm-continuous (if it were, by
      Stone's theorem the infinitesimal generator would be bounded). -/
  U_continuous : ∀ x : H, Continuous (fun g : G => U g x)
  covariance   : ∀ g a, (U g : H →L[K] H) * π a = π (σ g a) * (U g)
  U_commutes_F : ∀ g, (U g : H →L[K] H) * F = F * (U g)
\end{verbatim}

A graded variant adds \texttt{U\_commutes\_γ}. The analysis is unchanged
from the non-equivariant case: nothing is harder to prove; the new
fields are purely algebraic compatibility. Mathlib provides unitary
groups, continuous group homomorphisms, and strongly-continuous
representations (SOT continuity is exactly what Mathlib's
\texttt{ContinuousSMul\ G\ H} records, modulo unitarity), so this is a
small layer. We will introduce it as \texttt{Equivariant.lean} in Phase
4.5 \emph{after} the index pairing of Phase 2 has stabilized.

\subsubsection{Families --- heavy, defer but design for
it}\label{families-heavy-defer-but-design-for-it}

A ``family of Fredholm modules over a coefficient \(C^*\)-algebra
\(B\)'' is a Hilbert \(C^*\)-module \(\mathcal{E}\) over \(B\) together
with a bounded \(B\)-linear operator
\(F \in \mathcal{L}_B(\mathcal{E})\) satisfying the Fredholm-module
axioms with \texttt{IsCompactOperator} replaced by membership in the
\(B\)-compact ideal \(\mathcal{K}_B(\mathcal{E})\).

Mathlib status (as of \texttt{v4.29.0-rc8}):

\begin{itemize}
\tightlist
\item
  \texttt{Mathlib.Analysis.CStarAlgebra.Module.*} --- basic Hilbert
  \(C^*\)-module API exists.
\item
  \(\mathcal{L}_B(\mathcal{E})\) (adjointable operators) --- partial.
\item
  \(\mathcal{K}_B(\mathcal{E})\) (the closed two-sided ideal generated
  by \(\theta_{\xi,\eta}\)) --- not yet developed in usable form.
\item
  Regular self-adjoint unbounded operators on \(\mathcal{E}\) ---
  absent.
\end{itemize}

\textbf{Forward compatibility.} The Phase 1 definitions in
\texttt{Basic.lean} are written with \texttt{H\ :\ Type*} plus
typeclasses, \emph{not} with \(\mathbb{C}\) baked into the analysis
lemmas. To lift to families, the substitution will be:

{\def\LTcaptype{none} % do not increment counter
\begin{longtable}[]{@{}
  >{\raggedright\arraybackslash}p{(\linewidth - 2\tabcolsep) * \real{0.4940}}
  >{\raggedright\arraybackslash}p{(\linewidth - 2\tabcolsep) * \real{0.5060}}@{}}
\toprule\noalign{}
\begin{minipage}[b]{\linewidth}\raggedright
Single-space (current)
\end{minipage} & \begin{minipage}[b]{\linewidth}\raggedright
Families (Phase 6)
\end{minipage} \\
\midrule\noalign{}
\endhead
\bottomrule\noalign{}
\endlastfoot
\texttt{{[}InnerProductSpace\ K\ H{]}\ {[}CompleteSpace\ H{]}} &
\texttt{{[}HilbertCStarModule\ B\ E{]}} \\
\texttt{H\ →L{[}K{]}\ H} & \texttt{E\ →L{[}B{]}\ E} (adjointable
\(B\)-linear) \\
\texttt{IsCompactOperator\ (T\ :\ H\ →L{[}K{]}\ H)} &
\texttt{T\ ∈\ 𝒦\_B(E)} (the \(B\)-compact ideal) \\
Index lands in \texttt{ℤ} & Index lands in \texttt{K₀(B)} \\
\end{longtable}
}

The structural fields (one operator, three compactness axioms) and the
proof outline of the index pairing survive unchanged. The implementation
work for Phase 6 is the Mathlib gap --- not a redesign of our
structures.

\subsubsection{Combined: equivariant
families}\label{combined-equivariant-families}

The full generality is a \texttt{G}-equivariant Kasparov
\((\mathcal{A}, B)\)-bimodule, \(KK^G(\mathcal{A}, B)\). Out of scope
for this project. We list it for completeness: once Phases 4.5 and 6 are
in place, combining them is formally a product (both add commuting
layers of structure).

\subsection{3.9 Super-Bakry--Émery: encoding
sketch}\label{super-bakryuxe9mery-encoding-sketch}

The mathematical definitions are in §2.5; the encoding follows the same
mixin pattern as §3.8.

\begin{verbatim}
structure BakryEmeryFredholmModule (ρ : K) extends GradedFredholmModule K A H where
  /-  The carré du champ Γ(a, b) := ½ (δ(a)* δ(b) + δ(b*) δ(a*)*)
      with δ(a) := F * π a - π a * F. Returned as a bounded operator
      on H. -/
  Γ : A → A → H →L[K] H
  Γ_def : ∀ a b, Γ a b = ½ • ((δ a)* * δ b + (δ (star b)) * (δ (star a))*)
    where δ a := F * π a - π a * F
  /-  Iterated Γ via L = F². The factor convention matches §2.5. -/
  Γ₂ : A → A → H →L[K] H
  Γ₂_def : ∀ a b, Γ₂ a b = ½ • (L * Γ a b - Γ a (L_act b) - Γ (L_act a) b)
    where L := F * F
  /-  The CD(ρ, ∞) inequality: Γ₂(a, a) ≥ ρ • Γ(a, a) as a
      self-adjoint operator inequality. -/
  CD : ∀ a, Γ₂ a a - ρ • Γ a a ∈ {T : H →L[K] H | 0 ≤ T}
\end{verbatim}

Three places this differs from a naive translation:

\begin{enumerate}
\def\labelenumi{\arabic{enumi}.}
\tightlist
\item
  \textbf{\texttt{L\_act} is the dual action of \texttt{L} on
  \texttt{A}} --- not all \(L = F^2\) has an action on \(\mathcal{A}\)
  in the abstract setting. For commutative \(\mathcal{A}\) acting by
  multiplication, \(L_act(a)\) is defined by the unique element
  \(b \in \mathcal{A}\) with
  \(\pi(b) = L \pi(a) - \pi(a) L - [L, \pi(a)]_+\) when such \(b\)
  exists; otherwise the structure is formulated directly on the operator
  level (no action of \(L\) on \(\mathcal{A}\)). For the canonical
  manifold case \(L\) acts as the Laplacian on
  \(\mathcal{A}= C^\infty(M)\), so this is fine.
\item
  \textbf{Operator-positivity} \texttt{0\ ≤\ T} is via Mathlib's
  \texttt{IsSelfAdjoint\ T\ ∧\ ∀\ x,\ 0\ ≤\ ⟪x,\ T\ x⟫} (or the
  equivalent positive operator API in
  \texttt{Mathlib.Analysis.InnerProductSpace.Positive}).
\item
  \textbf{Equivalence theorem} (gradient-decay \(\Leftrightarrow\)
  \(CD\)): stated as
  \texttt{theorem\ BakryEmeryFredholmModule.gradient\_decay\_iff\_CD}.
  In the manifold case the proof goes via Lichnerowicz; abstractly via
  semigroup interpolation (Bakry's original argument).
\end{enumerate}

\subsubsection{\texorpdfstring{Bridge to
\texttt{graphops-qft}}{Bridge to graphops-qft}}\label{bridge-to-graphops-qft}

A finite \texttt{SusyGraphop} over \(\mathbb{R}\) is essentially a
finite-dimensional graded Fredholm module after complexification:

\begin{verbatim}
def SusyGraphop.toFredholmModule (S : SusyGraphop) :
    FredholmModule ℂ (Matrix (Fin S.n_plus) (Fin S.n_plus) ℂ)
                   ((Fin S.n_plus → ℂ) × (Fin S.n_minus → ℂ)) where
  π := Matrix.toLin'   -- diagonal-matrix action on the first factor
  F := blockMatrix 0 S.Q_minus.cmap S.Q_plus.cmap 0
       |>.normalize    -- F := (something)·Q with appropriate scaling
  ...
\end{verbatim}

with the grading
\(\gamma := \begin{pmatrix} 1 & 0 \\ 0 & -1 \end{pmatrix}\) on the
direct sum. The Witten index of \texttt{graphops-qft} then coincides
with our index pairing at \([p] = [1] \in K_0(\mathcal{A})\) (the unit
projection):
\(\langle [F], [1] \rangle = \dim \ker Q_+ - \dim \ker Q_- =
\mathrm{Witten}(S)\).

This bridge unifies the two projects under one abstraction: any theorem
about \texttt{BakryEmeryFredholmModule} specializes to a theorem about
\texttt{SusyGraphop} via \texttt{toFredholmModule}, and vice versa for
theorems specific to the finite case.

\section{4. Notation conventions}\label{notation-conventions}

Internal to the formalization we use:

\begin{itemize}
\tightlist
\item
  \texttt{K} --- base field (\texttt{ℝ} or \texttt{ℂ}).
\item
  \texttt{A} --- the algebra (in math: \(\mathcal{A}\)).
\item
  \texttt{H} --- the Hilbert space (in math: \(\mathcal{H}\)).
\item
  \texttt{π} --- the representation.
\item
  \texttt{F} --- the bounded operator (Fredholm-module form).
\item
  \texttt{D} --- the unbounded Dirac-type operator (spectral-triple
  form).
\item
  \texttt{γ} --- the chirality / grading.
\end{itemize}

For manifold-specific objects we add a subscript \texttt{M}:
\texttt{D\_M}, \texttt{γ\_M}, \texttt{H\_M\ =\ L\^{}2(M,\ S)}.

\section{5. Resolved decisions}\label{resolved-decisions}

The five questions in the original draft are answered below.

\begin{enumerate}
\def\labelenumi{\arabic{enumi}.}
\tightlist
\item
  \textbf{Unbounded operators --- yes for the spine (revised
  2026-08-20).} The core, resolvent criterion, canonical bounded
  commutator, and compact-resolvent package are implemented on
  \texttt{LinearPMap}. This is not finite summability. The bounded
  transform and general chiral pairing remain open.
\item
  \textbf{Algebra category --- \texttt{StarRing\ ℂ\ A} +
  \texttt{Algebra\ ℂ\ A}, not \texttt{CStarAlgebra}.} Demanding
  \(C^*\)-completeness would exclude \(C^\infty(M)\) (which is Fréchet,
  not \(C^*\)) and break the canonical manifold construction. Connes'
  framework was specifically designed over dense, holomorphically closed
  pre-\(C^*\)-algebras for this reason.
\item
  \textbf{Index encoding --- distinguish two current invariants.}
  \texttt{Fredholm.index} applies to bounded maps after an
  \texttt{IsFredholm} proof. \texttt{gradedKernelIndex} applies to the
  full even Dirac kernel and is not yet proved equal to a chiral
  Fredholm index. The trace formula
  \(\mathop{\mathrm{ind}}= \mathop{\mathrm{Tr}}(\gamma e^{-tD^2})\) is
  beautiful but requires heat semigroups, trace-class estimates, and the
  McKean--Singer asymptotic expansion --- none of which is in Mathlib.
  The algebraic definition is exact, native, and far easier to
  manipulate.
\item
  \textbf{Examples first --- yes; swap Phases 3 and 4.} The \(S^1\)
  Dirac triple is essentially a diagonal operator on
  \(\ell^2(\mathbb{Z})\); it tests the unbounded core, while the separate
  shift example tests the bounded Fredholm API. The chiral pairing
  remains open.
\item
  \textbf{Equivariance --- defer to Phase 4.5.} Baking
  \(G\)-equivariance into the Phase 1 definitions would complicate
  typeclass inference and force every non-equivariant theorem to carry
  an unused \texttt{{[}TopologicalGroup\ G{]}} instance. Lean's
  \texttt{extends} makes adding the equivariance layer cheap when we get
  to it.
\end{enumerate}

\section{6. Continuation: noncommutative metric geometry from the
algebraic
core}\label{continuation-noncommutative-metric-geometry-from-the-algebraic-core}

\emph{The phases above are scoped to definitions, the index pairing, the
spectral distance, and concrete examples. This section sketches the
theorems that \textbf{emerge automatically} from the same algebraic data
once both §2.5 (Bakry--Émery \(\Gamma_2\)) and §2.7 (Connes' spectral
distance) are in place. None of these is a deliverable of the current
plan; together they are the natural research-level continuation,
encompassing modern noncommutative metric geometry.}

The unifying observation: every result in this section follows from
operator inequalities on \(\Gamma\), \(\Gamma_2\), \([D, \pi(a)]\), and
the heat semigroup \(P_t = e^{-tD^2}\). No manifold tools --- no
geodesics, volume forms, Jacobi fields, or Levi--Civita connection ---
are required. Each theorem therefore specializes simultaneously to (i)
the canonical Dirac triple of a smooth manifold, (ii) the finite
\texttt{SusyGraphop} of a graph, and (iii) any further instance of the
abstract framework.

\subsection{6.1 Optimal transport (the Wasserstein-1
metric)}\label{optimal-transport-the-wasserstein-1-metric}

For commutative \(\mathcal{A}= C(X)\) with \(X\) compact, states are
probability measures on \(X\). The Connes distance evaluated on states
\(\phi_\mu, \phi_\nu\) corresponding to measures \(\mu, \nu\) satisfies,
by \textbf{Kantorovich--Rubinstein duality}: \[
d_D(\phi_\mu, \phi_\nu) \;=\; W_1(\mu, \nu)
\;=\; \inf_{\pi \in \Pi(\mu, \nu)} \int d_g(x, y)\,d\pi(x, y),
\] the Wasserstein-1 (earth-mover) distance. The Lipschitz condition
\(\|[D, \pi(a)]\| \leq 1\) is the dual Kantorovich condition. For the
graph case (finite \(\mathcal{A}\), Dirac masses), the same formula
reduces to the linear-programming form of \(W_1\) and recovers the
shortest-path metric. \textbf{One \texttt{sSup} definition absorbs the
entire optimal-transport formalism.}

\subsection{6.2 Entropic curvature and heat-flow contraction
(Lott--Sturm--Villani)}\label{entropic-curvature-and-heat-flow-contraction-lottsturmvillani}

The synthetic Ricci-curvature lower bound \(CD(\rho, \infty)\) is
\emph{equivalent} to the heat flow being a strict exponential
contraction on \((\mathcal{S}(\mathcal{A}), d_D)\): \[
d_D\bigl(\phi \circ P_t,\;\psi \circ P_t\bigr) \;\leq\; e^{-\rho t}\;d_D(\phi, \psi)
\qquad \forall\, t \geq 0,\ \phi, \psi \in \mathcal{S}(\mathcal{A}).
\] This is the von Renesse--Sturm characterization, lifted to the
noncommutative setting by Lott, Erbar--Maas, Carlen--Maas, and
Wirth--Zhang. \textbf{Heat distributions exponentially collapse together
at rate \(\rho\)} --- a pure metric-geometric reformulation of the
gradient-decay form of §2.5, useful as an \emph{alternative}
axiomatization of \(CD(\rho, \infty)\) that avoids operator-positivity.

\subsection{6.3 Poincaré inequality and spectral
gap}\label{poincaruxe9-inequality-and-spectral-gap}

Integrating \(\Gamma_2(a, a) \geq \rho\,\Gamma(a, a)\) along the heat
semigroup yields the noncommutative \textbf{Poincaré inequality}: for
any mean-zero state, the variance is controlled by the Dirichlet energy,
\[
\mathrm{Var}_\phi(a) \;\leq\; \tfrac{1}{\rho}\,\phi\bigl(\Gamma(a, a)\bigr) \qquad (\rho > 0).
\] Equivalently, the first non-zero eigenvalue of \(L = D^2\) satisfies
\(\lambda_1 \geq \rho\) (Lichnerowicz spectral gap for manifolds; mixing
rate for finite graphs / quantum Markov semigroups). The proof is a
two-line algebraic manipulation of \(P_t\) and requires no manifold
infrastructure --- strikingly different from the classical Lichnerowicz
proof, which uses the Bochner formula and integration by parts on the
sphere/manifold. In the discrete case it specializes to bounds on
random-walk convergence rates on finite graphs.

\subsection{6.4 Log-Sobolev inequality and
hypercontractivity}\label{log-sobolev-inequality-and-hypercontractivity}

A strictly positive curvature bound \(CD(\rho, \infty)\), \(\rho > 0\),
in fact implies the much stronger \textbf{logarithmic Sobolev
inequality} \[
\mathrm{Ent}_\phi(a^2) \;\leq\; \tfrac{2}{\rho}\,\phi\bigl(\Gamma(a, a)\bigr),
\] which by Gross's theorem is equivalent to
\textbf{hypercontractivity}: the heat semigroup \(P_t\) maps
\(L^2 \to L^p\) for \(p > 2\) in finite time, with explicit time
constants. This is the original purpose of the Bakry--Émery \(\Gamma_2\)
calculus (1985) and is the technical backbone of statistical-mechanics
phase transitions and (in the quantum-Markov-semigroup setting of
Junge--Mei--Parcet, Carlen--Maas) the sharpest known bounds on
noncommutative mixing. \textbf{All three inequalities --- Poincaré,
log-Sobolev, hypercontractivity --- fall out of the same operator
inequality \(\Gamma_2 \geq \rho \Gamma\).}

\subsection{6.5 Quantum Bonnet--Myers (diameter
bounds)}\label{quantum-bonnetmyers-diameter-bounds}

Upgrade the curvature condition to the \textbf{dimensional}
\(CD(\rho, n)\) inequality by adding a \(\tfrac{1}{n}(La)^2\)
correction: \[
\Gamma_2(a, a) \;\geq\; \rho\,\Gamma(a, a) \;+\; \tfrac{1}{n}\,(La)^2.
\] Then (Rieffel; Connes; Bakry--Qian) the \textbf{Connes--Wasserstein
diameter} of the state space is finite, \[
\mathrm{diam}\bigl(\mathcal{S}(\mathcal{A}), d_D\bigr) \;\leq\; \pi\sqrt{\tfrac{n}{\rho}}.
\] This is the \textbf{noncommutative Bonnet--Myers theorem}: positive
Ricci curvature forces compactness with a quantitative diameter bound
--- a purely operator-algebraic statement of the classical Riemannian
result, working uniformly for manifolds, graphs, and abstract
\(C^*\)-algebras. The classical proof uses Jacobi field ODEs along
geodesics; the algebraic proof uses three lines of \(\Gamma_2\)
calculus.

\subsection{6.6 Quantum Gromov--Hausdorff convergence (continuum
limits)}\label{quantum-gromovhausdorff-convergence-continuum-limits}

The \textbf{Lipschitz seminorm} \(L(a) := \|[D, \pi(a)]\|\) on a
spectral triple is the data of a \textbf{compact quantum metric space}
in Rieffel's sense (1999, 2004). The \textbf{quantum Gromov--Hausdorff
distance} \(d_{\mathrm{QGH}}\) defines a metric on the proper class of
all compact quantum metric spaces. Statement: \[
\bigl(\mathcal{A}_n, L_n\bigr) \xrightarrow{\;d_{\mathrm{QGH}} \to 0\;} \bigl(\mathcal{A}, L\bigr)
\] makes ``the lattice graphop converges to the smooth manifold'' into a
precise mathematical statement. For our Phase 3.5 \(T^2\) test: the
square-lattice \texttt{SusyGraphop} at mesh \(\varepsilon\) converges in
quantum Gromov--Hausdorff to the smooth Dirac
\texttt{GradedFredholmModule} of \(T^2\) as \(\varepsilon \to 0\). This
is the rigorous form of the \textbf{continuum-limit thesis} that
underwrites both this project and {[}\texttt{graphops-qft}{]}.

\subsection{6.7 What ties it together}\label{what-ties-it-together}

Each section above is a theorem about \(\Gamma\), \(\Gamma_2\),
\([D, \pi(a)]\), and \(P_t = e^{-tD^2}\). Each classical proof relies on
manifold-specific tools (geodesics, Jacobi fields, Bochner integration
by parts, volume forms, the Levi--Civita connection); the
algebraic-operator proofs replace all of these with operator
inequalities on the spectral triple. Once Phases 4.75 (super-BE) and 2.5
(spectral distance) are in place, \textbf{the same Lean theorem proves
each result simultaneously for smooth manifolds, finite graphs, and
arbitrary \texttt{BakryEmeryFredholmModule}s} --- including fractal
spaces, quantum groups, and the C*-algebraic deformations of
\texttt{graphops-qft}.

This is the long-term payoff of the synthesis: a uniform
metric-geometric calculus, formalized algebraically, that specializes to
the classical results without re-proof.

\section{7. References}\label{references}

\begin{enumerate}
\def\labelenumi{\arabic{enumi}.}
\tightlist
\item
  \textbf{Connes}, \emph{Noncommutative Geometry}, Academic Press, 1994.
  Chapter VI (definition + index pairing); Chapter IV §2 (local index
  formula with Moscovici).
\item
  \textbf{Gracia-Bondía, Várilly, Figueroa}, \emph{Elements of
  Noncommutative Geometry}, Birkhäuser, 2001. Chapters 9--11 (Dirac
  operator, canonical triple, axiomatics).
\item
  \textbf{Higson, Roe}, \emph{Analytic K-Homology}, OUP, 2000. Chapter 8
  (Fredholm modules, index pairing); Chapter 10 (Dirac).
\item
  \textbf{Kasparov}, ``The operator \(K\)-functor and extensions of
  \(C^*\)-algebras'', \emph{Izv. Akad. Nauk SSSR Ser. Mat.} 44 (1980),
  571--636 (English transl.: \emph{Math. USSR-Izv.} 16, 513--572);
  ``Equivariant \(KK\)-theory and the Novikov conjecture'',
  \emph{Invent. Math.} 91 (1988), 147--201. Foundational for the
  families and equivariant generalizations of §2.4 / §3.8.
\item
  \textbf{Atiyah, Singer}, ``The index of elliptic operators I, III'',
  \emph{Ann. of Math.} 87 (1968), 484--530 and 546--604; ``IV'' (1971),
  119--138. The classical index theorem and its families generalization;
  the target that Phases 4 and 6 specialize to in the manifold case. The
  equivariant fixed-point theorem of Phase 4.5 is from
  Atiyah--Segal--Singer (same series, paper II). 5a. \textbf{Atiyah,
  Patodi, Singer}, ``Spectral asymmetry and Riemannian geometry I, II,
  III'', \emph{Math. Proc. Camb. Phil. Soc.} 77 (1975) 43--69, 78 (1975)
  405--432, 79 (1976) 71--99. The index theorem for manifolds with
  boundary, the eta-invariant, and the APS non-local boundary conditions
  of §2.6's ``Manifolds with boundary'' out-of-scope note.
\item
  \textbf{Lawson, Michelsohn}, \emph{Spin Geometry}, Princeton
  University Press,

  \begin{enumerate}
  \def\labelenumii{\arabic{enumii}.}
  \setcounter{enumii}{1988}
  \tightlist
  \item
    Chapters I--II for Clifford algebras, \(\mathrm{Spin}(n)\), spinor
    bundles and Dirac operators --- the bundle-theoretic input to Phase
    4a.
  \end{enumerate}
\item
  \textbf{Bakry, Émery}, ``Diffusions hypercontractives'',
  \emph{Séminaire de Probabilités XIX}, Springer LNM 1123 (1985),
  177--206. The original \(\Gamma_2\) calculus and curvature-dimension
  condition referenced in §2.5 / Phase 4.75.
\item
  \textbf{Cipriani, Sauvageot}, ``Derivations as square roots of
  Dirichlet forms'', \emph{J. Funct. Anal.} 201 (2003), 78--120. The
  noncommutative carré du champ from a spectral triple --- the technical
  basis for the
  \(\Gamma(a, b) = \tfrac{1}{2}(\delta(a)^*\delta(b) + \dots)\) formula.
\item
  \textbf{Cushing, Kamtue, Liu, Münch, Peyerimhoff, Snodgrass},
  ``Bakry--Émery curvature sharpness and curvature flow in finite
  weighted graphs'', arXiv:2204.10064 (2022). Rigorous convergence
  theory for the BE flow on finite graphs, providing test cases for the
  \texttt{graphops-qft} bridge.
\item
  \textbf{\texttt{graphops-qft}} (related internal project). The
  finite-dimensional / graph-theoretic instantiation of the same
  \texttt{SusyGraphop\ =\ GradedFredholmModule} structure, with super-BE
  curvature flows and the spectral-triple connection worked out at the
  discrete level.
\item
  \textbf{Witten}, ``Supersymmetry and Morse theory'', \emph{J.
  Differential Geom.} 17 (1982), 661--692. The deformation
  \(d_t = e^{-tf} d e^{tf}\) and the heuristic Morse-complex picture
  referenced in §2.6 / Phase 3.7.
\item
  \textbf{Helffer, Sjöstrand}, ``Puits multiples en mécanique
  semi-classique IV: étude du complexe de Witten'', \emph{Comm. Partial
  Differential Equations} 10 (1985), 245--340. The rigorous
  semiclassical analysis of the \(t \to \infty\) limit of the Witten
  Laplacian, recovering Morse inequalities and the Morse complex.
\item
  \textbf{Connes}, ``Compact metric spaces, Fredholm modules, and
  hyperfiniteness'', \emph{Ergodic Theory Dynam. Systems} 9 (1989),
  207--220. The spectral-distance formula of §2.7.
\item
  \textbf{Rieffel}, ``Metrics on state spaces'', \emph{Doc. Math.} 4
  (1999), 559--600. The Lipschitz-seminorm /
  compact-quantum-metric-space framework --- when the spectral distance
  metrizes the weak-\(*\) topology on \(\mathcal{S}(\mathcal{A})\).
\end{enumerate}

\subsubsection{For §2.8 (cyclic
cohomology)}\label{for-2.8-cyclic-cohomology}

14a. \textbf{Connes}, ``Noncommutative differential geometry'',
\emph{Publ. Math. IHÉS} 62 (1985), 41--144. The original cyclic
cohomology paper and the \((b, B)\)-cocycle definition of the
Chern--Connes character of a Fredholm module. 14b. \textbf{Loday},
\emph{Cyclic Homology}, Springer Grundlehren 301, 2nd ed. (1998).
Textbook reference for Hochschild and cyclic (co)homology, the
\((b, B)\)-bicomplex, and periodicity. 14c. \textbf{Jaffe, Lesniewski,
Osterwalder}, ``Quantum K-theory I: The Chern character'', \emph{Comm.
Math. Phys.} 118 (1988), 1--14. The JLO entire-cyclic cocycle
representing \(\mathrm{ch}(D)\) in the unbounded picture. 14d.
\textbf{Hochschild, Kostant, Rosenberg}, ``Differential forms on regular
affine algebras'', \emph{Trans. Amer. Math. Soc.} 102 (1962), 383--408.
The original HKR theorem identifying Hochschild cohomology with
differential forms; the algebraic input to Connes' iso
\(HP^*(C^\infty(M)) \cong H^*_{\mathrm{dR}}(M)\).

\subsubsection{For §2.9 (noncommutative
instances)}\label{for-2.9-noncommutative-instances}

14e. \textbf{Connes}, ``\(C^*\)-algèbres et géométrie différentielle'',
\emph{C. R. Acad. Sci. Paris Sér. A-B} 290 (1980), A599--A604. The
original construction of the spectral triple on the noncommutative torus
\(A_\theta\), including the isospectral property. 14f. \textbf{Rieffel},
``\(C^*\)-algebras associated with irrational rotations'', \emph{Pacific
J. Math.} 93 (1981), 415--429. Construction of the Powers--Rieffel
projection \(p_\theta\) realizing \(\tau(p_\theta) = \theta\); the
K-theory of \(A_\theta\). 14g. \textbf{Pimsner, Voiculescu}, ``Exact
sequences for \(K\)-groups and Ext-groups of certain cross-product
\(C^*\)-algebras'', \emph{J. Operator Theory} 4 (1980), 93--118. The
Pimsner--Voiculescu exact sequence computing
\(K_0(A_\theta) \cong \mathbb{Z}+ \mathbb{Z}\) and the range of the
trace \(\tau(K_0(A_\theta)) = \mathbb{Z}+ \theta\mathbb{Z}\). 14h.
\textbf{Connes, Landi}, ``Noncommutative manifolds, the instanton
algebra and isospectral deformations'', \emph{Comm. Math. Phys.} 221
(2001), 141--159. The isospectral-deformation construction, generalizing
the NC torus to a broad class of noncommutative manifolds with the same
Dirac spectrum as a classical model.

\subsubsection{For §6 (continuation
directions)}\label{for-6-continuation-directions}

\begin{enumerate}
\def\labelenumi{\arabic{enumi}.}
\setcounter{enumi}{14}
\tightlist
\item
  \textbf{Villani}, \emph{Optimal Transport: Old and New}, Springer
  Grundlehren 338 (2009). The reference for Wasserstein distance,
  Kantorovich--Rubinstein duality (§6.1), and the Lott--Sturm--Villani
  synthetic-Ricci framework (§6.2).
\item
  \textbf{Sturm}, ``On the geometry of metric measure spaces I, II'',
  \emph{Acta Math.} 196 (2006), 65--177. The metric-measure-space
  formulation of \(CD(K, N)\) via entropy convexity in Wasserstein
  space.
\item
  \textbf{Lott, Villani}, ``Ricci curvature for metric-measure spaces
  via optimal transport'', \emph{Ann. of Math.} 169 (2009), 903--991.
  The independent (and complementary) synthetic formulation; jointly
  with {[}16{]} this is ``LSV theory''.
\item
  \textbf{Bakry, Gentil, Ledoux}, \emph{Analysis and Geometry of Markov
  Diffusion Operators}, Springer Grundlehren 348 (2014). The
  comprehensive reference for \(\Gamma_2\)-calculus, Poincaré, LSI,
  hypercontractivity, and the diffusion analogue of Bonnet--Myers
  (§§6.3--6.5).
\item
  \textbf{Rieffel}, ``Gromov--Hausdorff distance for quantum metric
  spaces'', \emph{Mem. Amer. Math. Soc.} 168 (2004), no. 796. Defines
  \(d_{\mathrm{QGH}}\) on the class of compact quantum metric spaces
  (§6.6).
\item
  \textbf{Carlen, Maas}, ``Gradient flow and entropy inequalities for
  quantum Markov semigroups with detailed balance'', \emph{J. Funct.
  Anal.} 273 (2017), 1810--1869. Entropic Ricci curvature for quantum
  Markov semigroups --- the noncommutative form of the
  Lott--Sturm--Villani contraction (§6.2).
\item
  \textbf{Wirth, Zhang}, ``Complete gradient estimates of quantum Markov
  semigroups'', \emph{Comm. Math. Phys.} 387 (2021), 761--791. Sharpest
  known \(\Gamma_2\)-style bounds in the noncommutative setting.
\end{enumerate}

\end{document}
